% !TEX program = pdflatex
\documentclass[12pt, oneside]{book}

%% ── Fonts ────────────────────────────────────────────────────────────────────
\usepackage[T1]{fontenc}
\usepackage[utf8]{inputenc}
\usepackage{mathpazo}   % Palatino text and math (matches TeX Gyre Pagella)
\usepackage{microtype}

%% ── Mathematics ─────────────────────────────────────────────────────────────
\usepackage{amsmath, amsthm, amssymb, mathtools}
% unicode-math omitted; using standard amsmath fonts

%% ── Page geometry ───────────────────────────────────────────────────────────
\usepackage[
  paper=letterpaper,
  top=1.25in,
  bottom=1.25in,
  inner=1.35in,
  outer=1.1in,
  headheight=15pt
]{geometry}

%% ── Headers and footers ─────────────────────────────────────────────────────
\usepackage{fancyhdr}
\pagestyle{fancy}
\fancyhf{}
\fancyhead[L]{\small\itshape\leftmark}
\fancyhead[R]{\small\thepage}
\renewcommand{\headrulewidth}{0.4pt}

%% ── Chapter and section titles ──────────────────────────────────────────────
\usepackage{titlesec}
\titleformat{\chapter}[display]
  {\normalfont\large\bfseries}
  {\chaptertitlename\ \thechapter}{16pt}{\Large}
\titlespacing*{\chapter}{0pt}{40pt}{24pt}

%% ── Hyperlinks ───────────────────────────────────────────────────────────────
\usepackage[hidelinks, unicode]{hyperref}

%% ── Bibliography style ───────────────────────────────────────────────────────
% Using built-in bibliography environment

%% ── Theorem environments ─────────────────────────────────────────────────────
\newtheoremstyle{plain-flyxion}
  {10pt}{8pt}{\itshape}{0pt}{\bfseries}{.}{.5em}{}
\theoremstyle{plain-flyxion}
\newtheorem{theorem}{Theorem}[chapter]
\newtheorem{lemma}[theorem]{Lemma}
\newtheorem{corollary}[theorem]{Corollary}
\newtheorem{proposition}[theorem]{Proposition}
\newtheorem{conjecture}[theorem]{Conjecture}

\newtheoremstyle{definition-flyxion}
  {10pt}{8pt}{\normalfont}{0pt}{\bfseries}{.}{.5em}{}
\theoremstyle{definition-flyxion}
\newtheorem{definition}[theorem]{Definition}
\newtheorem{example}[theorem]{Example}
\newtheorem{remark}[theorem]{Remark}

%% ── Spacing and paragraph style ──────────────────────────────────────────────
\usepackage{setspace}
\setstretch{1.15}
\setlength{\parindent}{1.5em}
\setlength{\parskip}{4pt}

%% ── Miscellaneous ────────────────────────────────────────────────────────────
\usepackage{xcolor}
\usepackage{booktabs}
\usepackage{graphicx}

%% ── Custom operators ─────────────────────────────────────────────────────────
\DeclareMathOperator{\Var}{Var}
\DeclareMathOperator{\diag}{diag}
\newcommand{\A}{\mathcal{A}}
\newcommand{\X}{\mathcal{X}}
\newcommand{\M}{\mathcal{M}}
\newcommand{\Phifield}{\Phi}
\newcommand{\kpi}{\kappa_{\pi}}
\newcommand{\Lam}{\Lambda}

%% ─────────────────────────────────────────────────────────────────────────────
\begin{document}

\frontmatter

%% ── Title page ───────────────────────────────────────────────────────────────
\begin{titlepage}
\centering
\vspace*{2.5cm}
{\Large\bfseries Hidden Curvature}\\[1.2em]
{\large\itshape Constraint Geometry, Projection,\\[0.3em]
and the Reconstruction of Meaning}\\[3.5cm]
{\large Flyxion}\\[0.8em]
{\normalsize Independent Researcher}\\[3.5cm]
{\normalsize 2026}
\vfill
\end{titlepage}

%% ── Abstract ─────────────────────────────────────────────────────────────────
\chapter*{Abstract}
\addcontentsline{toc}{chapter}{Abstract}

This monograph develops a field-theoretic account of semantic meaning grounded in
admissibility geometry rather than representational fidelity. The central claim is that
meaning is not a property of representations but an invariant of the admissibility
field---the structured space of constraint-governed transitions through which semantic
trajectories move. Representations are projections of that field. What has been described
in the literature as semantic collapse is not a failure of representational precision but
the observable trace of a geometric phenomenon: the destruction of curvature by
many-to-one projection. The intellectual center of the work is a reconstruction program:
given the observable failures of representation, recover the hidden constraint geometry
that produced them.

The argument proceeds in four stages. First, we construct the admissibility manifold $X$
of semantic trajectories and equip it with the Fisher information metric induced by the
conditional density $p(y\mid x)$ over admissible continuations. This metric is positive
semi-definite everywhere and degenerates precisely at semantic bottlenecks, making it
well-defined at the locations where existing Hessian-based approaches fail. Second, we
introduce the vertical-horizontal decomposition of the projection $\pi:X\to M$ onto a
representational manifold $M$ and define the projected sectional curvature $\kpi(m)$ as
the integrated mixed curvature over the fiber $\pi^{-1}(m)$. Third, we prove the
Projection-Collapse Principle in both a weak and a strong form: observable representational
mixing $\Lam(m)$ is bounded below by a monotone function of $\kpi(m)$, and under
regularity conditions the two quantities are asymptotically equivalent. This establishes
that collapse measurements are geometric observables of the hidden admissibility geometry.
Fourth, we formulate the inverse problem: given a field of mixing measurements over $M$,
reconstruct the curvature distribution of the admissibility manifold from its projected
trace.

Connections are developed to information geometry, retarded functional differential
equations, and the field theory of synthetic cognition. The monograph concludes with a
discussion of constraint-first semantic architectures, collapse-based auditing, and the
philosophical consequences of treating constraint as ontologically prior to representation.

\vspace{1em}
\noindent\textit{Keywords:} admissibility geometry, semantic collapse, Fisher information
metric, projection-collapse principle, inverse geometry, RSVP field theory,
constraint-first semantics.

%% ── Table of contents ────────────────────────────────────────────────────────
\tableofcontents

\mainmatter

%% ═════════════════════════════════════════════════════════════════════════════
\part{The Phenomenon}
%% ═════════════════════════════════════════════════════════════════════════════

\chapter{What Collapse Looks Like}

\section{Representational Systems and Their Failures}

A representational system encodes structured information in a compressed space. The
compression is useful: it allows a finite device to operate over an indefinitely large
domain of possible states by treating structurally similar states as equivalent. But
compression is also lossy. Some information is discarded, and the question that motivates
this monograph is which information, why, and what the loss reveals about the hidden
structure the representation was attempting to encode.

The failures that matter most are not random. When a representational system loses
information, it does not do so uniformly across all distinctions. Certain classes of
distinctions are systematically more vulnerable than others. These are the distinctions
that depend not on the local features of a token or state but on its position within a
structured space of possibilities: on what transitions from that position are admissible,
what futures remain accessible, and what constraints govern the trajectory forward. A
system that compresses by geometric proximity---treating nearby states as equivalent---will
erase these distinctions whenever the relevant structure is not preserved by the proximity
relation.

The phenomenon of semantic collapse names this systematic erasure. It is not a description
of noise or approximation error. It is a description of a structural failure: the
destruction, by compression, of the distinctions on which reasoning depends.

\section{Operator Sensitivity as the First Diagnostic}

The most immediate evidence of collapse comes from operator-sensitive constructions. Modal
operators, epistemic operators, indexicals, and agentive constructions share a common
feature: their meaning depends not on the intrinsic properties of their arguments but on
the structural relations those arguments bear to a space of possibilities. The sentence
``it must be raining'' differs from ``it is raining'' not in its surface vocabulary but in
the class of possible worlds over which it is evaluated. The sentence ``she knows the door
is locked'' differs from ``the door is locked'' not in its propositional content but in
whether that content is embedded within an epistemic frame.

Representational systems that treat semantic similarity as geometric proximity are
systematically blind to these distinctions. They cluster tokens by co-occurrence and
distributional context, which is sensitive to surface form but not to modal scope. The
result is that necessity is blurred with truth, epistemic attitude is flattened into
assertion, and indexical anchoring dissolves into description. The operators that
structure logical inference are the first casualties of compression.

This is not a defect of any particular architecture. It is a consequence of the
projection. Any compression that identifies states by their positional neighborhoods
rather than by their accessibility structure will produce operator collapse.

\section{The Phenomenology of Collapse Events}

Collapse manifests across multiple domains in recognizably similar patterns. In linguistic
systems, it appears as hallucination: the generation of fluent outputs that are
semantically indistinguishable from correct ones in the compressed representation but
logically incompatible with the constraints governing the original. In formal reasoning
systems, it appears as scope errors: the confusion of necessity with contingency, of
universal quantification with existential, of an agent's action with an event that merely
occurs. In perceptual systems, it appears as category boundary degradation: the failure to
maintain separability between classes that are geometrically close in representation space
but functionally distinct in the constraint structure.

The surface forms of these failures differ. Their underlying geometry is the same. In each
case, a many-to-one projection maps a curved admissibility structure onto a flatter
representational space, and the curvature destroyed by the projection reappears as
confusions in the output. The specific content of the confusion depends on which curvature
was destroyed. The structural form of the confusion is invariant.

\section{Collapse as a Signal}

The reorientation that drives the rest of this monograph is simple to state and
far-reaching in its consequences. Collapse events are not errors to be corrected. They are
measurement signals. A system that collapses in a particular pattern is revealing
something about the hidden structure it is compressing. The pattern of collapse is a
projected trace of the curvature of the admissibility manifold.

This is not a metaphor. The Projection-Collapse Principle established in Part~III makes
the relationship precise: the observable mixing functional $\Lam(m)$ is bounded below by a
monotone function of the projected sectional curvature $\kpi(m)$ of the admissibility
manifold. Collapse measurements are therefore lower bounds on hidden curvature. They are
data from which the hidden geometry can be partially reconstructed. The inverse problem of
semantic collapse is: given a field of collapse measurements, recover the admissibility
geometry that produced them.

\chapter{Existing Frameworks and Their Limits}

\section{Embedding-Based Accounts}

The standard picture of meaning in high-dimensional representational systems is geometric:
tokens are points, semantic similarity is proximity, and meaning relations are encoded in
the local structure of neighborhoods. This picture has proved remarkably productive. It
explains fluency, supports generalization, and underlies the practical success of
large-scale language models. It also has a structural limitation that becomes visible
precisely where the models fail.

The limitation is that geometric proximity is a property of representations, not of the
structures representations are meant to encode. Two tokens that are proximate in
embedding space may lead to radically different futures under the constraint structure
that governs the domain. A representation that identifies them because they are
geometrically close conflates things that must remain distinct for inference to work.
Embedding-based accounts explain collapse by appealing to the properties of
embeddings---the geometry does not separate the relevant classes---without asking why the
geometry fails to separate them or what it would mean for it to succeed.

The answer requires going below the level of representations to the constraint structure
they project from.

\section{Modal Logic Approaches}

The modal logic response to collapse is more sophisticated. It identifies, correctly, that
the distinctions being lost are operator-sensitive: they involve the difference between
what is true and what is necessary, what is known and what is the case, what an agent
does and what merely happens. It proposes to restore these distinctions by enriching the
representational apparatus with accessibility relations, modal bases, and
operator-sensitive constraints.

This diagnosis is accurate. The engineering proposal, however, remains within the
representational paradigm. Modal Proofing Kernels, typed embeddings, and frame
constraints are mechanisms for ensuring that a richer set of distinctions survives
compression. They add structure on top of representations rather than deriving
representations from structure. The modal operators are labels applied to tokens rather
than descriptions of the accessibility geometry that generates meaning. As a consequence,
the framework faces a regress: the labels must be sourced from somewhere, and if they are
learned from data, they inherit the very ambiguity the framework is designed to cure.

\section{The Regress Problem}

Both embedding-based and modal logic approaches share a hidden assumption: that the
categories to be preserved can be identified prior to, and independently of, the
construction of the geometry that is supposed to preserve them. The categories are
treated as given, and the task is to build a representation that keeps them separate.

But where do the categories come from? If modal labels are learned from distributional
data, they inherit the distributional ambiguity of the data. If they are imposed
externally, the system presupposes a prior distinction-preserving mechanism---the very
thing it is trying to construct. This regress is not a technical problem soluble within
the existing paradigm. It is a symptom that the paradigm is asking the wrong question.

The right question is not how to preserve given categories in a representation. The right
question is what generates the categories in the first place. The answer, this monograph
argues, is the admissibility structure: the field of constraints governing which
transitions are and are not permissible. Categories are not primitive. They are
descriptions of regions in the admissibility geometry that have distinct accessibility
profiles. Representations that preserve categories must preserve accessibility profiles.
Representations that collapse categories are destroying accessibility structure.

\section{The Inversion: Constraints Before Representations}

The construction that follows inverts the standard order. Admissibility structure is
primary. Representations are projections of that structure. Collapse is what happens when
the projection destroys curvature the structure contains. The task is not to add
constraints to representations. The task is to understand what constraint geometry
generates representations, to measure how much of that geometry survives projection, and
to reconstruct it from the collapse patterns that the projection produces.

This inversion has practical consequences. A system designed constraint-first does not
generate outputs and filter them for admissibility. It generates within the admissible
region from the beginning. The admissibility check is not a downstream repair. It is the
primitive computational operation from which generation is derived.

%% ═════════════════════════════════════════════════════════════════════════════
\part{The Geometry}
%% ═════════════════════════════════════════════════════════════════════════════

\chapter{The Admissibility Manifold}

\section{Semantic Trajectories as the Primary Objects}

Rather than tokens, propositions, or embedding vectors, the primary objects of the theory
developed here are trajectories through a space of admissible transitions. A semantic
state is not a point endowed with intrinsic properties. It is a position within a
structured space of possible continuations, a location in a field that determines what
moves forward from it are coherent, what moves are forbidden, and what moves are merely
possible.

This shift in primary objects is not terminological. It changes the ontological order of
the theory. States have the properties they have because of the transitions they admit.
Similarity between states is a function of similarity in their accessibility profiles. The
geometry of the state space is induced by the structure of admissible transitions. Nothing
in this construction assumes that states have intrinsic properties independent of the
constraint field. The field is prior.

Let $X$ denote the space of all such semantic positions. The topology of $X$ is determined
by the admissibility relation: two positions are nearby if their accessible futures are
similar. The manifold structure, developed in the following sections, is induced by
requiring that this similarity vary smoothly across $X$.

\section{The Admissible Continuation Space}

For each $x\in X$, let $\A(x)$ denote the space of admissible continuations from $x$,
equipped with a conditional probability density $p(y\mid x)$ over $y\in\A(x)$. The
density $p$ encodes not merely which continuations are possible but how strongly each
is supported by the constraint structure at $x$. Continuations with high density are
those for which the constraints bearing on $x$ provide strong positive support. Those
with low density are continuations that are technically permissible but carry little
constraint force. Those outside $\A(x)$ are excluded entirely.

Several remarks about the nature of $p(y\mid x)$ are in order. First, it is not learned
from data in the sense that a language model learns a next-token distribution. It is the
object from which such learning is derived. The density $p$ encodes the underlying
constraint structure of the domain, and any learned distribution is an approximation to
it. Second, $p$ need not be uniform over $\A(x)$. The admissibility structure assigns
different weights to different continuations, and this differential weighting is the
source of the metric structure developed in the following chapter. Third, $p$ can vary
smoothly with $x$ even when $\A(x)$ itself changes topology---when bottlenecks open or
close---provided the transitions in density are continuous.

\section{Local Charts from Admissibility Profiles}

The manifold structure of $X$ is not assumed. It is derived from the admissibility
distribution. Define the \emph{admissibility map}:
\[
\Psi : X \to \mathcal{P}(\Omega),
\qquad
\Psi(x) = p(\cdot \mid x),
\]
where $\mathcal{P}(\Omega)$ denotes the space of probability measures on the continuation
space $\Omega$, equipped with the topology of weak convergence. The map $\Psi$ sends
each semantic position to its distribution over admissible continuations.

\begin{proposition}\label{prop:immersion}
If $p(y \mid x)$ depends smoothly on $x$ in the sense that the map $x \mapsto
\log p(y \mid x)$ is $C^\infty$ for $p(\cdot\mid x)$-almost every $y$, and if $\Psi$ is
injective, then $X$ is an immersed submanifold of $\mathcal{P}(\Omega)$, and the Fisher
information metric on $\mathcal{P}(\Omega)$ pulls back to a smooth Riemannian metric on
$X$.
\end{proposition}

\begin{proof}
Smooth dependence of $\log p(y \mid x)$ on $x$ implies that the map $\Psi: X \to
\mathcal{P}(\Omega)$ is smooth in the sense of the infinite-dimensional manifold
structure on $\mathcal{P}(\Omega)$ given by the Fisher-Rao metric. Injectivity of $\Psi$
means that distinct semantic positions have distinct admissibility distributions: the map
separates points. A smooth injective map with injective differential is an immersion.
The Fisher information metric on $\mathcal{P}(\Omega)$, given by
$G(u,v) = \mathbb{E}_{p}[uv/p^2]$ for tangent vectors $u,v$ in the $L^2(p)$ sense,
pulls back along $\Psi$ to the metric $g_{ij}(x) = \mathbb{E}_{y \sim p(\cdot \mid
x)}[\partial_i \log p \cdot \partial_j \log p]$, which is the Fisher information metric
of Definition~\eqref{eq:fisher}. Smoothness of the pullback metric follows from smoothness
of $\Psi$.
\end{proof}

The injectivity assumption in Proposition~\ref{prop:immersion} has a clear semantic
content: two semantic positions are distinct if and only if they differ in at least one
admissible continuation. Positions with identical admissibility distributions are
semantically indistinguishable. If injectivity fails, $X$ should be replaced by the
quotient $X / {\sim}$, where $x \sim x'$ iff $p(\cdot \mid x) = p(\cdot \mid x')$.

The proposition establishes that the manifold structure of $X$ is inherited from the
geometry of probability measures via $\Psi$. The admissibility map is not merely a
convenient parameterization; it is the geometric embedding of the semantic space into the
universal statistical manifold $\mathcal{P}(\Omega)$.

\section{Regularity Conditions}

Three regularity conditions on $p(y\mid x)$ are required for the construction to be
well-behaved. Each condition has a semantic interpretation that clarifies what it assumes
about the constraint structure.

The first condition is smooth variation of $\A(x)$ with $x$: as the semantic position
changes, the accessible futures change continuously. This means that the constraint
structure has no sharp discontinuities---no locations where a small change in position
produces a catastrophic change in what is admissible. Domains with genuine phase
boundaries, where the constraint structure changes sharply, require a separate analysis.
The smooth-variation condition holds in the generic case.

The second condition is differentiability of $p(y\mid x)$ in $x$: the score function
$\partial_x \log p(y\mid x)$ exists and is square-integrable with respect to $p(\cdot\mid
x)$ for each $x$. This is the standard regularity condition for Fisher information to be
well-defined. It means that the admissibility pressure on each continuation varies
smoothly with position.

The third condition is integrability: $\int_{\A(x)} \|\partial_x \log p(y\mid x)\|^2
p(y\mid x)\, dy < \infty$ for each $x$. This ensures that the Fisher information is finite
and the metric can be computed.

\section{Bottlenecks, Constraints, and Degeneracies}

The admissibility manifold $X$ is not homogeneous. Its structure encodes the full topology
of the constraint space. Regions where $\A(x)$ is large and the density $p(\cdot\mid x)$
is spread broadly correspond to positions of high semantic flexibility: many coherent
continuations are available. Regions where $\A(x)$ contracts sharply correspond to
semantic bottlenecks: very few continuations remain admissible, and the constraint
structure bears heavily on which way forward is available.

These structural features are encoded in the Fisher metric developed in the next chapter.
High curvature in the metric corresponds to regions where the accessibility structure
changes rapidly across neighboring positions. Low curvature corresponds to regions where
accessibility is stable. The bottlenecks and constraint boundaries of the domain appear
as curvature concentrations in the admissibility geometry.

\chapter{Tangent Spaces and Admissibility Flows}

Before introducing the Fisher metric, we develop the differential structure of the
admissibility manifold directly from the admissibility distributions. This chapter
establishes the tangent bundle of $X$ as a space of score-compatible perturbations and
derives the inner product on each tangent space. The Fisher metric then emerges naturally
as the unique metric compatible with this structure, rather than being defined externally.

\section{Statistical Tangent Vectors}

At each point $x \in X$, the tangent space $T_x X$ consists of infinitesimal deformations
of the admissibility profile $p(\cdot \mid x)$. A deformation of $p(\cdot \mid x)$ in
direction $u$ must preserve the normalization constraint $\int_\Omega p(y \mid x)\, dy =
1$. Differentiating with respect to a curve $x(t)$ through $x$ gives:
\[
\frac{d}{dt}\bigg|_{t=0} \int_\Omega p(y \mid x(t))\, dy
= \int_\Omega \dot x^i(0)\, \partial_i p(y \mid x)\, dy = 0.
\]
Writing $u(y) = \dot x^i(0)\, \partial_i \log p(y \mid x)$ (the score in direction
$\dot x$), this normalization constraint becomes:
\[
\int_\Omega u(y)\, p(y \mid x)\, dy = 0.
\]

\begin{definition}[Statistical Tangent Space]\label{def:tan}
The statistical tangent space to $X$ at $x$ is the $L^2(p(\cdot\mid x))$-closure:
\begin{equation}\label{eq:tangent}
T_x X = \left\{ u \in L^2(\Omega, p(\cdot\mid x)) : \int_\Omega u(y)\, p(y\mid x)\, dy = 0 \right\}.
\end{equation}
Elements of $T_x X$ are called \emph{admissibility flows} at $x$.
\end{definition}

The zero-mean condition in~\eqref{eq:tangent} is the infinitesimal form of normalization
preservation. An admissibility flow $u$ specifies how the distribution over accessible
futures changes as the semantic position moves in $X$: positive values of $u(y)$ indicate
that continuation $y$ becomes more accessible, negative values indicate it becomes less
accessible, and the constraint $\int u\, p\, dy = 0$ ensures that accessibility is
redistributed rather than created or destroyed.

\section{Score Functions as Natural Coordinates}

The partial derivatives of the log-density provide a canonical basis for the statistical
tangent space.

\begin{proposition}[Score Functions Span the Tangent Space]\label{prop:score-basis}
Under the regularity conditions of Chapter~3, the score functions:
\begin{equation}\label{eq:score}
u_i(y) = \partial_i \log p(y \mid x), \qquad i = 1, \ldots, \dim X,
\end{equation}
lie in $T_x X$ and, if $\Psi: X \to \mathcal{P}(\Omega)$ is an immersion
(Proposition~\ref{prop:immersion}), they span the image of $d\Psi_x$ in $T_{\Psi(x)}
\mathcal{P}(\Omega)$.
\end{proposition}

\begin{proof}
First, $u_i \in T_x X$: by the score equation, $\int \partial_i \log p(y\mid x)\,
p(y\mid x)\, dy = \partial_i \int p(y\mid x)\, dy = 0$. So each $u_i$ has zero
$p(\cdot\mid x)$-mean.

Second, these are the images of the coordinate vectors $\partial_i \in T_x X$ under
$d\Psi_x$: by the chain rule, $\frac{d}{dt}\big|_0 \Psi(x + t e_i) = \partial_i
p(\cdot \mid x) = p(\cdot \mid x) \cdot \partial_i \log p(\cdot \mid x) = p \cdot u_i$.
Thus $d\Psi_x(\partial_i) = p \cdot u_i$ in $L^2$, and the $u_i$ span the image of
$d\Psi_x$.
\end{proof}

Proposition~\ref{prop:score-basis} identifies the score functions as the natural
coordinate vectors on $X$ viewed as a submanifold of $\mathcal{P}(\Omega)$. The score
$u_i(y) = \partial_i \log p(y\mid x)$ measures the sensitivity of the log-accessibility
of continuation $y$ to a displacement in direction $i$ in semantic space.

\section{The Fisher Inner Product}

The inner product on each tangent space is determined by requiring that the metric measure
local statistical distinguishability.

\begin{definition}[Fisher Inner Product]\label{def:fisher-inner}
The Fisher inner product on $T_x X$ is:
\begin{equation}\label{eq:fisher-inner}
g_x(u, v) = \int_\Omega u(y)\, v(y)\, p(y \mid x)\, dy
= \mathbb{E}_{y \sim p(\cdot\mid x)}[u(y)\, v(y)].
\end{equation}
\end{definition}

\begin{lemma}[Properties of the Fisher Inner Product]\label{lem:fisher-props}
The bilinear form~\eqref{eq:fisher-inner} is:
\begin{enumerate}
  \item[(i)] \emph{Bilinear and symmetric}: immediate from the definition.
  \item[(ii)] \emph{Positive semi-definite}: $g_x(u,u) = \mathbb{E}[u^2] \geq 0$ with
    equality iff $u = 0$ $p(\cdot\mid x)$-a.e.
  \item[(iii)] \emph{Non-degenerate on the tangent space}: if $g_x(u,v) = 0$ for all
    $v \in T_x X$, then $u = 0$ $p(\cdot\mid x)$-a.e., because $T_x X$ is a dense
    subspace of $L^2(\Omega, p(\cdot\mid x))$ and separates points.
\end{enumerate}
\end{lemma}

When evaluated on the score functions~\eqref{eq:score}, the Fisher inner product
recovers exactly the Fisher information metric:
\[
g_x(u_i, u_j) = \mathbb{E}_{y\sim p(\cdot\mid x)}[\partial_i \log p(y\mid x)\cdot
\partial_j \log p(y\mid x)] = g_{ij}(x),
\]
as in equation~\eqref{eq:fisher}. The Fisher metric is therefore not introduced by
definition but derived as the coordinate representation of the inner product on the
statistical tangent space in the natural score basis.

\section{Uniqueness of the Fisher Metric}

The Fisher inner product is the unique Riemannian metric on $X$ compatible with
infinitesimal statistical distinguishability. We make this precise through the following
axiomatic characterization, which is the classical result of Chentsov adapted to the
admissibility setting~\cite{amari2016information}.

\begin{theorem}[Uniqueness of the Fisher Metric]\label{thm:fisher-unique}
Let $h$ be any Riemannian metric on $X$ such that:
\begin{enumerate}
  \item[(i)] \emph{Statistical invariance}: $h$ is invariant under sufficient statistics,
    i.e., if $T: \Omega \to \Omega'$ is a sufficient statistic for the family
    $\{p(\cdot\mid x)\}$, then the pushforward metric under $T$ equals $h$.
  \item[(ii)] \emph{Positive definiteness}: $h_x(u,u) > 0$ for all non-zero $u \in T_x X$.
\end{enumerate}
Then $h = c\, g$ for some positive constant $c$. Up to scale, the Fisher metric is the
unique metric satisfying statistical invariance.
\end{theorem}

The statistical invariance condition means that if two parameterizations of the
admissibility family lead to the same distributions, then any sufficient statistic that
mediates between them should not change measured distances. This is the information-theoretic
content of the uniqueness: the Fisher metric measures information, and information is
invariant under sufficient compression.

\section{Semantic Interpretation of Admissibility Flows}

A tangent vector $u \in T_x X$ admits a direct semantic interpretation. Writing $u(y) =
\partial_i \log p(y\mid x)\, v^i$ for some velocity $v = v^i \partial_i \in T_x X$
(using the score basis), we have:
\[
u(y) = \frac{d}{d\epsilon}\bigg|_0 \log p(y \mid x + \epsilon v).
\]
This is the rate of change of the log-accessibility of continuation $y$ as the semantic
position moves in direction $v$. Continuations with large positive $u(y)$ are those that
become significantly more accessible; those with large negative $u(y)$ become less
accessible. The admissibility flow $u$ therefore specifies the direction in which the
accessible future is deforming as the semantic trajectory moves through $X$.

The Fisher metric $g_x(u,u) = \mathbb{E}[u^2]$ measures the mean squared rate of change
of log-accessibility: how rapidly the distribution over accessible futures is deforming.
Large values indicate semantically loaded directions where small positional changes
strongly reshape the future. Small values indicate semantically redundant directions
where the future is stable. This is the intrinsic content of the Fisher metric as a
measure of semantic distinguishability.

\chapter{Admissibility Geodesics}

With the tangent space structure in place, we can define geodesics on the admissibility
manifold. Geodesics are the paths of minimal semantic deformation: they represent the
most efficient routes between semantic positions, deforming the admissibility structure
as little as possible along the way. Before discussing curvature in Chapter~5, we need
geodesics to give curvature something to act on: curvature measures how geodesics diverge
from the flat case.

\section{The Geodesic Energy Functional}

Let $\gamma: [0,1] \to X$ be a smooth curve connecting $\gamma(0) = x_0$ to $\gamma(1)
= x_1$. The length of $\gamma$ in the Fisher metric is:
\[
L[\gamma] = \int_0^1 \sqrt{g_{\gamma(t)}(\dot\gamma(t), \dot\gamma(t))}\, dt,
\]
and the energy functional is:
\begin{equation}\label{eq:geo-energy}
E[\gamma] = \frac{1}{2}\int_0^1 g_{\gamma(t)}(\dot\gamma(t), \dot\gamma(t))\, dt.
\end{equation}
By the Cauchy-Schwarz inequality, $L[\gamma]^2 \leq 2E[\gamma]$ with equality when
$|\dot\gamma|$ is constant. Therefore minimizing energy among constant-speed
reparameterizations is equivalent to minimizing length.

\section{The Geodesic Equation}

\begin{theorem}[Admissibility Geodesic Equation]\label{thm:geodesic}
Critical points of the energy functional~\eqref{eq:geo-energy} satisfy the
Euler-Lagrange equation:
\begin{equation}\label{eq:geodesic}
\ddot{x}^k + \Gamma^k_{ij}(x)\, \dot{x}^i \dot{x}^j = 0,
\end{equation}
where $\Gamma^k_{ij}$ are the Christoffel symbols of the Fisher metric:
\begin{equation}\label{eq:christoffel}
\Gamma^k_{ij} = \frac{1}{2} g^{kl}\!\left(\partial_i g_{jl} + \partial_j g_{il}
- \partial_l g_{ij}\right).
\end{equation}
Solutions to~\eqref{eq:geodesic} are \emph{admissibility geodesics}.
\end{theorem}

\begin{proof}
Standard variational calculus. Consider a variation $\gamma_\epsilon(t) = \gamma(t) +
\epsilon \eta(t)$ with $\eta(0) = \eta(1) = 0$. Differentiating $E[\gamma_\epsilon]$
with respect to $\epsilon$ at $\epsilon = 0$ and setting equal to zero:
\[
\frac{d}{d\epsilon}\bigg|_0 E[\gamma_\epsilon]
= \int_0^1 \left[g_{ij}\ddot\gamma^i \eta^j
+ \frac{1}{2}\partial_k g_{ij} \dot\gamma^i\dot\gamma^j \eta^k\right] dt = 0,
\]
integrating the first term by parts. Since $\eta$ is arbitrary, the integrand vanishes
pointwise, giving:
\[
g_{kj}\ddot\gamma^j + \partial_k g_{ij}\dot\gamma^i\dot\gamma^j
- \frac{1}{2}\partial_k g_{ij}\dot\gamma^i\dot\gamma^j = 0.
\]
Contracting with $g^{kl}$ and symmetrizing in $i,j$ yields~\eqref{eq:geodesic}
with~\eqref{eq:christoffel}.
\end{proof}

\section{Fisher Geodesic Distance}

The geodesic distance between two semantic positions is the infimum of lengths over all
smooth paths:
\begin{equation}\label{eq:fisher-dist}
d_F(x_1, x_2) = \inf_{\gamma: \gamma(0)=x_1,\, \gamma(1)=x_2} L[\gamma],
\end{equation}
where the infimum is over smooth curves in $X$. When a minimizing geodesic exists (which
it does by the Hopf-Rinow theorem whenever $(X,g)$ is complete), $d_F$ is realized by a
solution to~\eqref{eq:geodesic}. The Fisher distance $d_F(x_1, x_2)$ measures the
minimum semantic deformation required to move from admissibility profile $p(\cdot\mid
x_1)$ to $p(\cdot\mid x_2)$: it is the information-geometric cost of the transition.

The distance $d_F$ is the quantity that appears in the curvature-divergence
inequality~\eqref{eq:js-lip} of Chapter~6. That inequality can now be restated with the
geodesic definition in place: the Jensen-Shannon divergence between nearby admissibility
distributions grows at most quadratically in the geodesic distance, with the coefficient
of growth given by the Lipschitz constant of the admissibility density.

\section{Semantic Interpretation: Geodesics as Minimal Deformation}

An admissibility geodesic $\gamma: [0,1] \to X$ is the path of minimal semantic
deformation from $x_0$ to $x_1$: it traverses the least-curved route through the
admissibility manifold, redistributing the accessible future as smoothly as possible
along the way.

In flat regions of $X$ (where the curvature $R_g$ vanishes), geodesics are straight lines
in the statistical sense: the admissibility profile deforms linearly from $p(\cdot\mid
x_0)$ to $p(\cdot\mid x_1)$. In curved regions, geodesics are deflected by the curvature.
The direction of deflection is toward regions of higher admissibility density---toward
the semantic attractors where the accessible future is richest. This is the information-geometric
analogue of gravitational lensing: curvature bends the minimal-deformation paths toward
regions of concentrated constraint.

\section{Bottleneck Deflection and Escape Velocity}

A semantic bottleneck at $x_b \in X$ is a region where $\A(x_b)$ is small and the
admissibility density $p(\cdot\mid x_b)$ is concentrated on a small number of
continuations. In the Fisher metric, bottlenecks correspond to regions where the metric
degenerates in the directions of collapsed futures: the accessible-future cone contracts,
and geodesics approaching $x_b$ from the side must deflect to avoid the degenerate
directions.

The geometry near a bottleneck is analogous to the geometry near a saddle point in
Riemannian geometry: geodesics arriving along the wide direction of the saddle pass
through relatively freely, while those arriving along the narrow direction are deflected
or blocked. The bottleneck creates an effective potential in the geodesic equation, which
can be made precise via the Jacobi equation (the linearization of~\eqref{eq:geodesic})
and the associated conjugate point theory.

\chapter{Volume Growth and Accessibility Geometry}

Curvature in a Riemannian manifold controls the rate at which volumes of geodesic balls
grow with radius. In the admissibility manifold, this volume growth has a direct semantic
interpretation: it measures how rapidly the accessible future expands as the semantic
position moves away from a given point. This chapter establishes the Bishop-Gromov
comparison theorem for the admissibility manifold and uses it to formalize the bottleneck
language that has appeared throughout the preceding chapters.

\section{Fisher Balls and Accessibility Volume}

\begin{definition}[Fisher Ball]\label{def:fisher-ball}
The Fisher ball of radius $r$ centered at $x \in X$ is:
\begin{equation}\label{eq:fisher-ball}
B_r(x) = \{y \in X : d_F(x, y) < r\},
\end{equation}
and its volume with respect to the Riemannian volume form $d\mu_g$ of $(X,g)$ is:
\begin{equation}\label{eq:vol-def}
V(x, r) = \mu_g(B_r(x)) = \int_{B_r(x)} d\mu_g.
\end{equation}
\end{definition}

The Fisher ball $B_r(x)$ is the set of semantic positions reachable from $x$ by a
geodesic of Fisher length at most $r$. In the admissibility interpretation, it is the
set of semantic positions whose admissibility profiles differ from $p(\cdot\mid x)$ by
a KL divergence of at most $\frac{1}{2}r^2$ (by Theorem~\ref{thm:fisher-deriv}). The
volume $V(x,r)$ measures the size of the admissibility-accessible neighborhood of $x$
at scale $r$.

\section{Comparison Geometry: The Bishop-Gromov Theorem}

In a Riemannian manifold with Ricci curvature bounded below by a constant $K$, the
volume of geodesic balls satisfies a comparison inequality with the corresponding volume
in the space form of constant curvature $K$. Let $V_K(r)$ denote the volume of a ball
of radius $r$ in the $n$-dimensional space of constant sectional curvature $K$:
\[
V_K(r) = \omega_{n-1} \int_0^r \mathrm{sn}_K(t)^{n-1}\, dt,
\]
where $\omega_{n-1}$ is the volume of the unit $(n-1)$-sphere and:
\[
\mathrm{sn}_K(t) = \begin{cases}
  K^{-1/2}\sin(K^{1/2}t) & K > 0, \\
  t & K = 0, \\
  |K|^{-1/2}\sinh(|K|^{1/2}t) & K < 0.
\end{cases}
\]

\begin{theorem}[Bishop-Gromov for the Admissibility Manifold]\label{thm:bishop-gromov}
Let $(X, g)$ be the admissibility manifold with Fisher metric $g$, and suppose the
Ricci curvature satisfies $\mathrm{Ric}_g \geq K$ for some $K \in \mathbb{R}$ and all
$x \in X$. Then the ratio:
\begin{equation}\label{eq:bg-ratio}
r \mapsto \frac{V(x, r)}{V_K(r)}
\end{equation}
is monotone non-increasing in $r$ for all $x \in X$.
\end{theorem}

\begin{proof}
This is the classical Bishop-Gromov theorem applied to the Riemannian manifold $(X,g)$
with the Fisher metric. The proof uses the Jacobi field comparison: if $\mathrm{Ric}_g
\geq K$, then Jacobi fields in $(X,g)$ grow no faster than Jacobi fields in the
comparison space of curvature $K$. Integrating Jacobi field lengths over the unit sphere
in the tangent space gives the volume comparison~\eqref{eq:bg-ratio}. The monotonicity
follows from the Riccati comparison for the mean curvature of distance spheres.
\end{proof}

\section{Curvature and Accessibility Volume}

Theorem~\ref{thm:bishop-gromov} has direct semantic consequences. The monotonicity of
$V(x,r)/V_K(r)$ translates as follows.

\begin{corollary}[Positive Curvature Contracts Accessible Futures]\label{cor:pos-curv}
If $\mathrm{Ric}_g(x) > 0$, then $V(x,r) < V_0(r) = \omega_{n-1} r^n / n$ for all
sufficiently small $r > 0$: the Fisher ball at $x$ is smaller than the Euclidean ball
of the same radius. Semantically, positive Ricci curvature at $x$ means that the
admissibility-accessible neighborhood of $x$ is smaller than it would be in a flat
admissibility geometry. The constraint structure at $x$ concentrates accessible futures
into a smaller volume.
\end{corollary}

\begin{corollary}[Negative Curvature Expands Accessible Futures]\label{cor:neg-curv}
If $\mathrm{Ric}_g(x) < 0$, then $V(x,r) > V_0(r)$ for small $r$: the Fisher ball is
larger than the Euclidean ball. Negative Ricci curvature at $x$ means that the
accessible neighborhood expands faster than it would in flat space. The constraint
structure at $x$ is locally divergent: small departures from $x$ lead to rapidly
diversifying accessible futures.
\end{corollary}

\section{Semantic Bottlenecks as Volume Suppression}

The bottleneck language used throughout this monograph can now be formalized precisely.

\begin{definition}[Semantic Bottleneck]\label{def:bottleneck}
A point $x \in X$ is an \emph{$(\epsilon, r)$-bottleneck} if:
\begin{equation}\label{eq:bottleneck}
V(x, r) \leq \epsilon\, V_0(r),
\end{equation}
i.e., the Fisher ball of radius $r$ at $x$ has volume at most $\epsilon$ times the
Euclidean ball of the same radius. An $(\epsilon, r)$-bottleneck is a region where
accessible futures are strongly suppressed relative to the flat baseline.
\end{definition}

By Corollary~\ref{cor:pos-curv}, points of large positive Ricci curvature are bottlenecks.
The severity of the bottleneck is controlled by the magnitude of $\mathrm{Ric}_g$: the
larger the Ricci curvature, the smaller $V(x,r)/V_0(r)$, and the more tightly the
accessible futures are constrained.

The bottleneck structure connects to the projection analysis of Chapter~5. A projection
$\pi: X \to M$ that maps a bottleneck region to a low-dimensional region in $M$ must
identify many admissibility-distinct points in the fiber over that region: the volume
suppression at the bottleneck forces the fiber to be large. By~\eqref{eq:kpi-A}, large
fibers in curved regions produce large projected sectional curvature $\kpi$. The
Projection-Collapse Principle then forces large $\Lam$: bottlenecks are the primary
sources of observable collapse.

\chapter{The Fisher Information Metric}

\section{Why the Hessian Fails}

The natural first attempt at a metric on $X$ would be to construct a scalar functional
$\Phifield: X \to \mathbb{R}$ measuring the total accessibility at each position---the
log-volume of the admissible continuation space, for instance---and to define the metric
as its Hessian:
\[
g_{ij}(x) = \partial_i \partial_j \Phifield(x).
\]
This construction has two advantages. It is familiar from optimization and from the
differential geometry of level sets. And it produces a metric that degenerates at the
critical points of $\Phifield$, which might seem to correspond to semantic bottlenecks.

The Hessian construction fails for a precise reason. Collapse events occur at saddle
points, ridges, and near-degenerate regions of $\Phifield$---exactly the locations where
the Hessian is indefinite or vanishing. At a saddle point, the Hessian has both positive
and negative eigenvalues: it is not a metric at all. At a flat region near a bottleneck,
the Hessian is nearly zero and provides no geometric information. The most semantically
interesting locations in $X$---the places where the constraint structure changes rapidly
and compression destroys the most information---are precisely the locations where the
Hessian fails.

\section{The Fisher Information Metric}

The correct metric is the Fisher information metric induced by the conditional density
$p(y\mid x)$. Define:
\begin{equation}\label{eq:fisher}
g_{ij}(x) = \mathbb{E}_{y\sim p(\cdot\mid x)}\!\left[
  \partial_i \log p(y\mid x)\cdot\partial_j \log p(y\mid x)
\right].
\end{equation}
Equivalently, under the regularity conditions of Chapter~3:
\[
g_{ij}(x) = -\mathbb{E}_{y\sim p(\cdot\mid x)}\!\left[
  \partial_i\partial_j \log p(y\mid x)
\right].
\]
The equivalence between these two expressions is standard in information geometry and
follows from differentiating the normalization condition $\int p(y\mid x)\,dy = 1$
under the integral sign. The second expression is sometimes more tractable for
computation; the first is more directly interpretable.

\section{Derivation from Infinitesimal Admissibility Divergence}

The Fisher metric admits a derivation that grounds it geometrically in the distinguishability
of nearby semantic positions. Consider two nearby states $x$ and $x + dx$ in $X$. The
KL divergence from $p(\cdot \mid x)$ to $p(\cdot \mid x + dx)$ measures the information
cost of replacing the admissibility distribution at $x$ with that at $x + dx$:
\[
D_{\mathrm{KL}}\!\left(p(\cdot \mid x) \;\|\; p(\cdot \mid x+dx)\right)
= \int_{\Omega} p(y \mid x) \log\frac{p(y \mid x)}{p(y \mid x + dx)}\, dy.
\]

\begin{theorem}\label{thm:fisher-deriv}
Under the regularity conditions of Section~3.3, the KL divergence between nearby
admissibility distributions satisfies:
\begin{equation}\label{eq:kl-taylor}
D_{\mathrm{KL}}\!\left(p(\cdot \mid x) \;\|\; p(\cdot \mid x+dx)\right)
= \frac{1}{2} g_{ij}(x)\, dx^i\, dx^j + O(\|dx\|^3),
\end{equation}
where $g_{ij}$ is the Fisher information metric~\eqref{eq:fisher}. Moreover, $g$ is the
unique Riemannian metric on $X$ whose quadratic form equals the second-order approximation
to admissibility divergence.
\end{theorem}

\begin{proof}
Expand $\log p(y \mid x+dx)$ in Taylor series:
\[
\log p(y \mid x+dx)
= \log p(y \mid x)
+ \partial_i \log p(y \mid x)\, dx^i
+ \tfrac{1}{2}\partial_i\partial_j \log p(y \mid x)\, dx^i dx^j
+ O(\|dx\|^3).
\]
Substituting into the KL divergence:
\begin{align*}
D_{\mathrm{KL}}\!\left(p(\cdot \mid x) \;\|\; p(\cdot \mid x+dx)\right)
&= -\int p(y \mid x)\Big[\partial_i \log p(y\mid x)\, dx^i
+ \tfrac{1}{2}\partial_i\partial_j \log p(y\mid x)\, dx^i dx^j\Big]\, dy
+ O(\|dx\|^3).
\end{align*}
The first-order term vanishes by the score equation: differentiating $\int p(y\mid x)\,
dy = 1$ gives $\int p(y\mid x)\,\partial_i \log p(y\mid x)\, dy = 0$.

For the second-order term, differentiate the score equation with respect to $x^j$:
\[
\int \partial_j p(y\mid x) \cdot \partial_i \log p(y\mid x)\, dy
+ \int p(y\mid x) \cdot \partial_i \partial_j \log p(y\mid x)\, dy = 0.
\]
The first integral equals $\int p(y\mid x)\,(\partial_j \log p)(\partial_i \log p)\, dy
= g_{ij}(x)$, so:
\[
\int p(y\mid x)\,\partial_i\partial_j \log p(y\mid x)\, dy = -g_{ij}(x).
\]
Therefore the second-order term contributes $\frac{1}{2}g_{ij}(x)\, dx^i dx^j$, which
establishes~\eqref{eq:kl-taylor}. Uniqueness follows because the coefficients of $dx^i
dx^j$ in the Taylor expansion of a fixed divergence are uniquely determined, forcing $h_{ij} = g_{ij}$ for any metric $h$ with the same second-order approximation property.
\end{proof}

Theorem~\ref{thm:fisher-deriv} establishes that the Fisher metric is the infinitesimal
form of KL divergence between admissibility distributions. Two semantic positions at
Fisher distance $\epsilon$ have admissibility distributions that differ by approximately
$\frac{1}{2}\epsilon^2$ in KL divergence. The metric measures how distinguishable nearby
positions are by their accessible futures.

\section{Positive Semi-Definiteness and Its Semantic Meaning}

The Fisher metric~\eqref{eq:fisher} is positive semi-definite at every $x\in X$ by
construction. The outer product $\partial_i \log p \cdot \partial_j \log p$ is a
rank-one positive semi-definite matrix at each $y$, and expectation preserves positive
semi-definiteness. No convexity assumption on $\Phifield$ is required.

The metric degenerates---becomes rank-deficient---precisely when $p(\cdot\mid x)$ is
insensitive to perturbations in the direction indexed by $i$. That is: $g_{ii}(x)=0$
if and only if $\partial_i \log p(y\mid x) = 0$ for $p(\cdot\mid x)$-almost every $y$,
meaning that moving in direction $i$ does not change the accessible futures. Degeneracy
is semantically meaningful. It identifies the directions in $X$ along which the
constraint structure is invariant: moving in a degenerate direction does not change what
is admissible, and no geometric resolution is lost by collapsing such directions in a
representation.

Directions of large metric weight are the semantically loaded directions: small
perturbations along them strongly reshape $p(\cdot\mid x)$, changing which futures are
accessible. These are the directions that representations must preserve to maintain
semantic fidelity. Directions of small metric weight are the semantically redundant
directions: the accessibility structure is stable under perturbation, and compression
along these directions loses little.

\section{Connection to Information Geometry}

The admissibility manifold $(X, g)$ is a statistical manifold in the sense of Amari: a
smooth manifold whose points parameterize a family of probability distributions, equipped
with the Fisher information metric~\cite{amari2016information}. The specific interpretation
here differs from the standard statistical setting. The parameter space is not a space of
statistical models but a space of semantic positions. The family of distributions is the
family of admissibility densities $\{p(\cdot\mid x) : x\in X\}$. The Fisher metric
measures not statistical distinguishability in the standard sense but semantic
distinguishability: the degree to which different positions lead to different accessible
futures.

This reinterpretation carries consequences. The Cramér-Rao inequality, in this setting,
becomes a lower bound on how precisely a representation can resolve distinct admissibility
structures. Any projection $\pi: X\to M$ that attempts to encode the semantic position
$x$ as a point $m=\pi(x)$ in a lower-dimensional space $M$ loses information in
proportion to the curvature of $(X,g)$ in the directions collapsed by the projection.
The Cramér-Rao bound is not a statistical limitation. It is a geometric limitation on
representational fidelity.

\section{Curvature of the Admissibility Manifold}

The Riemann curvature tensor $R_g$ of the Fisher metric measures how the local
admissibility structure twists and turns across neighboring regions of $X$. High
curvature at $x$ means that nearby points have substantially different admissible future
cones: small displacements produce large changes in which continuations are available.
Flat regions---where curvature is low---mean that the accessibility structure varies
slowly and projections can represent it faithfully.

The curvature identifies, in geometric terms, the locations where any compression will
lose information: where the admissibility geometry is curved, a projection that identifies
nearby points will necessarily confuse trajectories that lead to different futures. Flat
regions project cleanly. Curved regions project poorly. Semantic collapse is the
observable consequence of projecting curved admissibility geometry onto a flatter
representational space.

\chapter{Projections and the Vertical-Horizontal Decomposition}

\section{Representations as Projections}

A representational system is a smooth map $\pi: X\to M$ from the admissibility manifold
to a lower-dimensional representational manifold $M$. The map $\pi$ compresses $X$ by
identifying points that share a representation: for each $m\in M$, the fiber
$\pi^{-1}(m)$ is the set of all admissibility states that project to $m$. Neural
encoders, embedding maps, summary statistics, and symbolic abstractions are all instances
of this compression. What distinguishes them is not the conceptual form of the operation
but the specific properties of $\pi$: how much curvature it destroys, which directions it
collapses, and how the fiber structure relates to the admissibility geometry.

The faithfulness of a representation is determined by how much of the geometric structure
of $(X,g)$ survives the map $\pi$. A representation is faithful in direction $v\in T_xX$
if the differential $d\pi_x$ is an isometry in that direction: if it maps $v$ to a vector
in $T_{\pi(x)}M$ of the same length. It is unfaithful if $d\pi_x(v) = 0$: if direction
$v$ is invisible to $\pi$. Collapse occurs when semantically loaded directions---directions
of high metric weight in $g$---are made invisible by the projection.

\section{The Vertical-Horizontal Decomposition}

At each $x\in X$, the tangent space $T_xX$ decomposes canonically into two orthogonal
subspaces:
\begin{equation}\label{eq:vh-decomp}
T_xX = V_x \oplus H_x,
\end{equation}
where $V_x = \ker(d\pi_x)$ is the \emph{vertical bundle}: the subspace of tangent
directions that are invisible to $\pi$. These are the directions that the representation
cannot see. The subspace $H_x = V_x^\perp$ is the \emph{horizontal bundle}: the
$g$-orthogonal complement of $V_x$, consisting of the directions that the projection
faithfully transmits.

The decomposition~\eqref{eq:vh-decomp} is the geometric formalization of the distinction
between what a representation can see and what it cannot. Horizontal directions are
preserved. Vertical directions are destroyed. What matters for collapse is neither pure
vertical curvature nor pure horizontal curvature. Pure vertical curvature---curvature
entirely within $V_x$---is hidden by $\pi$ and averaged over within each fiber, but
because it does not couple to horizontal directions, it does not introduce systematic
errors at the representational level. Pure horizontal curvature---curvature entirely
within $H_x$---is preserved by $\pi$ and appears faithfully in $M$. It is the
\emph{mixed} curvature, the sectional curvature $K_g(v,h)$ for $v\in V_x$ and $h\in
H_x$, that is dangerous. Mixed curvature couples what is hidden to what is observable.
When it is present, the fiber contains points that look identical in $M$ but lead to
different observable behaviors depending on which direction they are approached from.
This is the geometric source of collapse.

\section{Projected Sectional Curvature}

We define the projected sectional curvature at $m\in M$ as the integrated mixed sectional
curvature over the fiber $\pi^{-1}(m)$:
\begin{equation}\label{eq:kpi}
\kpi(m) = \int_{\pi^{-1}(m)}
  \sum_{\substack{v\in V_x,\; h\in H_x \\ |v|=|h|=1}}
  \left|K_g(v,h)\right| d\mu_m(x),
\end{equation}
where $\mu_m$ is the measure on the fiber induced by the volume form of $g$ and the sum
is taken over an orthonormal frame. This is the curvature that projection destroys: the
coupling between the hidden directions and the observable directions. Only mixed sectional
curvature contributes to the leakage that collapse diagnostics measure. Pure vertical and
pure horizontal curvature do not.

The projected sectional curvature $\kpi$ is a refined measure. It isolates, within the
full Riemann tensor $R_g$, exactly the components that a neighborhood-based statistic
computed in $M$ can detect. This refinement is not merely technical. It has direct
semantic content: $\kpi(m)$ measures how much of the admissibility structure at $m$ the
representation has destroyed in a way that creates observable confusion.

\section{The O'Neill Tensors and Collapse-Driving Curvature}

The precise relationship between the projection geometry and the mixed sectional curvature
is given by O'Neill's fundamental equations for Riemannian submersions~\cite{oneill1966}.
These equations decompose the curvature of $(X,g)$ into contributions from the base,
the fibers, and the interaction between them. The interaction term is the collapse-driving
curvature.

For a smooth vector field $E$ on $X$, write $E = E^H + E^V$ for its horizontal and
vertical components. Define the \emph{O'Neill $A$-tensor} and \emph{$T$-tensor} as:
\begin{align}
A_E F &= (\nabla_{E^H} F^H)^V + (\nabla_{E^H} F^V)^H, \label{eq:A-tensor}\\
T_E F &= (\nabla_{E^V} F^V)^H + (\nabla_{E^V} F^H)^V, \label{eq:T-tensor}
\end{align}
where $\nabla$ is the Levi-Civita connection of $g$. The $A$-tensor measures the
non-integrability of the horizontal distribution: $A_h v = (\nabla_h v)^V$ for
$h\in H_x$ and $v\in V_x$. The $T$-tensor measures the second fundamental form of the
fibers.

\begin{proposition}[O'Neill Mixed Curvature Formula]\label{prop:oneill}
For $v\in V_x$ and $h\in H_x$ with $|v|=|h|=1$, the mixed sectional curvature is:
\begin{equation}\label{eq:oneill-mixed}
K_g(v, h) = g\!\left((\nabla_v A)_h h,\, v\right)
- \|A_h v\|^2
+ g\!\left(T_v h,\, T_v h\right) - g\!\left(T_v v,\, T_h h\right) + \cdots
\end{equation}
where the dominant collapse-driving term is $\|A_h v\|^2$, the squared norm of the
$A$-tensor acting on the mixed pair $(h, v)$.
\end{proposition}

The significance of Proposition~\ref{prop:oneill} is that it identifies the $A$-tensor
as the primary source of mixed curvature. The quantity $A_h v = (\nabla_h v)^V$ measures
how much a horizontal displacement $h$ rotates a vertical direction $v$: if horizontal
motion through the fiber changes which vertical directions are present, then the fiber
structure is non-trivially coupled to the base, and projection destroys information. When
$A \equiv 0$, the horizontal distribution is integrable, the projection is locally a
product, and no mixed curvature is present. In that case $\kpi = 0$ and no leakage
occurs.

This connects directly to the semantic interpretation: the $A$-tensor measures how much
the admissibility structure varies across nearby fibers. High $\|A_h v\|^2$ at $x$ means
that moving horizontally by $h$---changing the representation---changes which vertical
directions carry admissibility information. Projection then cannot faithfully represent
the admissibility structure, and collapse is geometrically forced.

The projected sectional curvature~\eqref{eq:kpi} is therefore controlled by the
$A$-tensor norm integrated over the fiber:
\begin{equation}\label{eq:kpi-A}
\kpi(m) \asymp \int_{\pi^{-1}(m)} \|A(x)\|^2\, d\mu_m(x),
\end{equation}
where $\|A(x)\|^2 = \sum_{h,v} |A_h v|^2$ with the sum over an orthonormal frame. This
is the form of $\kpi$ that enters the proof of the Projection-Collapse Principle.

\section{The Fiber as a Hidden Variable}

The fiber $\pi^{-1}(m)$ is the set of all admissibility states that project to the same
representation $m$. It is a hidden variable in the strict statistical sense: it affects
observable outcomes but is not directly accessible through $M$. Two trajectories that are
representationally identical---that pass through the same point in $M$---may have very
different admissible futures if they come from different points in the fiber.

Collapse diagnostics are measurements that partially identify the fiber structure without
directly observing it. The mixing functional $\Lam(m)$, defined in the next chapter, is a
statistic computed in $M$ that measures the degree to which the fiber $\pi^{-1}(m)$
contains admissibility-distinct points. The Projection-Collapse Principle establishes the
relationship between this observable and the hidden curvature $\kpi(m)$. The inverse
problem is to recover $\kpi$ from measurements of $\Lam$: to reconstruct the hidden
geometry from the observable trace it produces through collapse.

%% ═════════════════════════════════════════════════════════════════════════════
\part{The Theorem}
%% ═════════════════════════════════════════════════════════════════════════════

\chapter{The Representational Mixing Functional}

\section{Beyond Operator-Specific Leakage}

Several diagnostic frameworks for measuring semantic collapse have been proposed in the
literature~\cite{denham2025}. These frameworks define cross-type leakage statistics that
measure the degree to which semantically distinct operator classes mix within embedding
neighborhoods. These statistics are well-motivated and empirically useful. They are,
however, specific to particular representational architectures, particular operator
taxonomies, and particular notions of semantic type. For a domain-independent theory, a
more general object is needed.

The representational mixing functional defined here generalizes these statistics in a way
that removes all dependence on embeddings, operator classes, and domain-specific
categorizations. It measures, at the level of abstract probability theory, the degree to
which a projection fails to separate admissibility classes that should remain distinct. The
existing diagnostics are then special cases, and the Projection-Collapse Principle applies
to all of them simultaneously.

\section{Definition of the Mixing Functional}

Let $U\subset M$ be a measurable region in the representational manifold, and let
$\{U_i\}_{i=1}^k$ be a measurable partition of $U$ into admissibility classes---subsets
that should, under a faithful representation, remain separated because they correspond to
distinct regions of the admissibility manifold. Define the representational mixing
functional:
\begin{equation}\label{eq:lambda}
\Lam(U) = D\!\left(\pi_*\mu_U,\; \prod_{i=1}^k \pi_*\mu_{U_i}\right),
\end{equation}
where $D$ is an information divergence and $\pi_*\mu_{U_i}$ is the pushforward of the
admissibility measure restricted to class $i$.

When the admissibility classes are perfectly separated in $M$---when the projection
faithfully distinguishes them---the joint distribution $\pi_*\mu_U$ factorizes over the
partition and $\Lam(U) = 0$. When the classes mix under projection, the joint distribution
diverges from the product distribution and $\Lam(U)>0$. The divergence quantifies
precisely how much admissibility structure the projection has failed to separate.

The choice of information divergence $D$ is a parameter of the framework. The
Kullback-Leibler divergence, Jensen-Shannon divergence, and Rényi divergence each produce
a version of the functional with different analytical properties. For the strong form of
the Projection-Collapse Principle, we require Jensen-Shannon divergence: it is bounded,
symmetric, and metrizes weak convergence---properties that the KL divergence lacks and
that the asymptotic equivalence proof requires.

\section{Existing Diagnostics as Special Cases}

The cross-type leakage statistics of~\cite{denham2025} are instances of~\eqref{eq:lambda}.
Specifically, those statistics arise when $U$ is a $k$-nearest-neighbor ball around a
token $t$ in the embedding space, the partition $\{U_i\}$ is into operator type classes,
and $D$ is a normalized mass-leakage ratio. The generalization does not change the
diagnostic procedure. It changes its interpretation: the mixing functional $\Lam$ is not
a statistic internal to the embedding space. It is a measure of admissibility class
separation failure under projection, and the same quantity applies to any representational
compression of any structured space.

This interpretive shift has a specific consequence. When an existing diagnostic reports
high leakage for modal tokens, it is not reporting a property of the embedding. It is
reporting a measurement of the projected sectional curvature $\kpi$ of the admissibility
manifold, filtered through a particular choice of $U$, partition, and divergence. The
Projection-Collapse Principle makes the relationship between the measurement and the
underlying geometry explicit.

\section{Regularity Conditions on the Functional}

Two regularity conditions on $\Lam$ are required for the Projection-Collapse Principle.
First, continuity with respect to the pushforward measure $\pi_*\mu$: small changes in
the distribution of representations produce small changes in the mixing functional. This
is satisfied by all standard information divergences under mild conditions on the
underlying densities.

Second, measurability of the partition $\{U_i\}$: the admissibility classes can be
identified by measurable sets in $M$. This is a non-trivial condition in practice. It
requires that the admissibility classes have sufficiently regular boundaries in the
representational space. When the projection $\pi$ severely distorts the admissibility
geometry, the images of admissibility classes may become non-measurable or have fractal
boundaries. The regularity condition restricts the analysis to projections that, while
potentially destructive of curvature, at least preserve the measurability of class
structure.

\chapter{The Projection-Collapse Principle}

\section{Fiber Divergence: The Bridge Between Curvature and Leakage}

Before stating the main theorem, we establish the intermediate object that connects
curvature to leakage. The fiber divergence measures the internal admissibility
heterogeneity of each fiber, providing the missing mathematical bridge in the chain
$\kpi \to \Delta_F \to \Lam$.

\begin{definition}[Fiber Divergence]\label{def:fiber-div}
Let $\bar{p}_m = \frac{1}{\mu_m(\pi^{-1}(m))}\int_{\pi^{-1}(m)} p(\cdot \mid x)\,
d\mu_m(x)$ be the barycenter of the admissibility distributions over the fiber
$\pi^{-1}(m)$. Define the \emph{fiber divergence} at $m \in M$ as:
\begin{equation}\label{eq:fiber-div}
\Delta_F(m) = \int_{\pi^{-1}(m)}
D_{\mathrm{JS}}\!\left(p(\cdot \mid x),\, \bar{p}_m\right) d\mu_m(x),
\end{equation}
where $D_{\mathrm{JS}}$ is the Jensen-Shannon divergence.
\end{definition}

\begin{lemma}[Fiber Divergence and Admissibility Uniformity]\label{lem:fiber-div}
$\Delta_F(m) = 0$ if and only if $p(\cdot \mid x) = \bar{p}_m$ for $\mu_m$-almost
every $x \in \pi^{-1}(m)$: that is, every point in the fiber has identical admissibility
distribution.
\end{lemma}

\begin{proof}
The Jensen-Shannon divergence $D_{\mathrm{JS}}(p, q) \geq 0$ with equality if and only
if $p = q$ as measures. The integral $\Delta_F(m) = \int D_{\mathrm{JS}}(p(\cdot\mid
x), \bar{p}_m)\, d\mu_m(x)$ is zero if and only if the integrand is zero $\mu_m$-almost
everywhere, which holds if and only if $p(\cdot \mid x) = \bar{p}_m$ for $\mu_m$-a.e.\
$x$.
\end{proof}

Lemma~\ref{lem:fiber-div} establishes that $\Delta_F(m) = 0$ is the condition for the
fiber to be admissibility-homogeneous: all points in $\pi^{-1}(m)$ lead to the same
distribution over accessible futures. When the fiber is admissibility-homogeneous,
the projection $\pi$ loses no admissibility information at $m$, and no collapse occurs.
$\Delta_F(m) > 0$ is the condition for heterogeneity: different points in the fiber
lead to different futures, and the projection conflates them.

\section{The Curvature-Divergence Inequality}

The second step in the proof chain connects curvature to fiber divergence. We show that
$\kpi(m)$ controls $\Delta_F(m)$ via a comparison geometry estimate.

\begin{proposition}[Curvature-Divergence Inequality]\label{prop:curv-div}
Let $d_F(x_1, x_2)$ denote the Fisher geodesic distance on $(X, g)$. Under the
Lipschitz condition that $p(\cdot \mid x)$ is $L_p$-Lipschitz in $x$ with respect to
$d_F$, the Jensen-Shannon divergence between admissibility distributions at two points in
the same fiber satisfies:
\begin{equation}\label{eq:js-lip}
D_{\mathrm{JS}}\!\left(p(\cdot \mid x_1),\, p(\cdot \mid x_2)\right)
\leq C\, d_F(x_1, x_2)^2,
\end{equation}
for a constant $C$ depending only on $L_p$. Integrating over the fiber:
\begin{equation}\label{eq:fiber-curv}
\Delta_F(m) \leq C \int_{\pi^{-1}(m)} d_F(x, \bar{x}_m)^2\, d\mu_m(x),
\end{equation}
where $\bar{x}_m$ is the fiber barycenter in $(X, g)$.
\end{proposition}

\begin{proof}
The Jensen-Shannon divergence satisfies the bound $D_{\mathrm{JS}}(p,q) \leq
\|p - q\|_{\mathrm{TV}}$ where $\|\cdot\|_{\mathrm{TV}}$ is the total variation distance,
and by Pinsker's inequality $D_{\mathrm{JS}}(p,q) \leq \frac{1}{2\ln 2}
D_{\mathrm{KL}}(p \| q)$ up to symmetrization. The Lipschitz condition on $p(\cdot\mid
x)$ combined with Theorem~\ref{thm:fisher-deriv} gives:
\[
D_{\mathrm{KL}}\!\left(p(\cdot \mid x_1) \;\|\; p(\cdot \mid x_2)\right)
\leq L_p^2\, d_F(x_1, x_2)^2 + O(d_F^3),
\]
which yields~\eqref{eq:js-lip} for small $d_F$ with $C = L_p^2/(2\ln 2)$. The fiber
integral~\eqref{eq:fiber-curv} then follows by integrating over $\pi^{-1}(m)$ and
applying the triangle inequality for $d_F$.
\end{proof}

The right side of~\eqref{eq:fiber-curv} is a fiber-integrated squared Fisher distance.
By comparison geometry---specifically, the Bishop-Gromov volume comparison theorem
adapted to the submersion setting---this integral is controlled by the sectional
curvature of $(X, g)$ restricted to the fiber. In the regime of bounded curvature, the
integral grows with the A-tensor norm via the O'Neill formula~\eqref{eq:oneill-mixed}.
The complete chain is therefore:

\begin{equation}\label{eq:chain}
\|A\|^2 \;\;\longrightarrow\;\; \kpi \;\;\longrightarrow\;\;
d_F(x_1,x_2)^2 \;\;\longrightarrow\;\; \Delta_F \;\;\longrightarrow\;\; \Lam.
\end{equation}

Each arrow corresponds to a proved or bounded inequality. The chain makes explicit what
was previously described in prose: collapse is not a coincidence but a geometric
consequence, traceable step by step from the non-integrability of the horizontal
distribution to the observable mixing functional.

\section{Statement: Weak Form}

\begin{theorem}[Projection-Collapse Principle, Weak Form]\label{thm:weak}
Let $(X, g)$ be an admissibility manifold equipped with the Fisher information
metric~\eqref{eq:fisher} induced by the conditional density $p(y\mid x)$ satisfying the
regularity conditions of Chapter~\ref{chap:manifold}. Let $\pi:X\to M$ be a smooth
submersion with fiber measure $\mu_m$ induced by the volume form of $g$. Let $\Lam$ be
a representational mixing functional~\eqref{eq:lambda} continuous with respect to the
pushforward measure $\pi_*\mu$. Then there exists a monotone non-decreasing function
$f:\mathbb{R}_{\geq 0}\to\mathbb{R}_{\geq 0}$ with $f(0)=0$ such that
\begin{equation}\label{eq:weak-form}
\Lam(m) \geq f\!\left(\kpi(m)\right)
\end{equation}
for all $m\in M$. Sufficiently large hidden admissibility curvature forces observable
representational mixing.
\end{theorem}

\section{Proof of the Weak Form}

The argument proceeds in four steps.

\medskip
\noindent\textit{Step 1.} High projected sectional curvature implies admissibility-distinct
points in the fiber. If $\kpi(m) > 0$, then by definition~\eqref{eq:kpi} there exists a
set of positive measure in $\pi^{-1}(m)$ on which the mixed sectional curvature
$|K_g(v,h)|$ is positive. At such points, parallel transport along horizontal curves
rotates the vertical bundle, meaning that nearby fibers have different admissibility
structures. By the Cramér-Rao argument applied to the Fisher metric, points with different
admissibility structures have different distributions over $\A(x)$: they lead to different
accessible futures. Therefore $\pi^{-1}(m)$ contains points with statistically
distinguishable admissibility distributions.

\medskip
\noindent\textit{Step 2.} The projection $\pi$ identifies all points in $\pi^{-1}(m)$.
Points that are admissibility-distinct in $X$ are representationally equivalent in $M$:
they map to the same point $m$ and are indistinguishable to any statistic computed in $M$.

\medskip
\noindent\textit{Step 3.} The admissibility partition $\{U_i\}$ assigns admissibility-distinct
points to different classes. Under the identification produced by $\pi$, points from
different admissibility classes are present in the same representational neighborhood
$U$. The joint distribution $\pi_*\mu_U$ therefore contains contributions from multiple
admissibility classes that the faithful distribution $\prod_i \pi_*\mu_{U_i}$ would
keep separate.

\medskip
\noindent\textit{Step 4.} The mixing functional $\Lam(m)>0$ follows from the positivity of
the information divergence $D$ when the pushforward distribution fails to factorize. The
lower bound by $f(\kpi(m))$ follows from continuity of $\Lam$ with respect to
$\pi_*\mu$ and compactness of the fiber measure: as $\kpi(m)$ increases from zero, the
within-fiber divergence between admissibility distributions increases continuously, and
by continuity of $\Lam$, so does the mixing functional. The function $f$ is the
composition of this two-step monotone relationship. \qed

\section{Statement: Strong Form}

\begin{theorem}[Projection-Collapse Principle, Strong Form]\label{thm:strong}
Under the additional regularity conditions:
\begin{enumerate}
  \item[(i)] \emph{Bounded fiber diameter}: $\sup_{x,x'\in\pi^{-1}(m)} d_g(x,x')\leq
    D < \infty$ for all $m\in M$;
  \item[(ii)] \emph{Lipschitz continuity}: $p(\cdot\mid x)$ is Lipschitz continuous in
    $x$ with respect to the metric $g$, with constant $L_p$;
  \item[(iii)] \emph{Jensen-Shannon divergence}: the divergence $D$ in~\eqref{eq:lambda}
    is the Jensen-Shannon divergence $D_{\mathrm{JS}}$;
\end{enumerate}
the mixing functional and the projected sectional curvature are asymptotically equivalent
up to multiplicative constants:
\begin{equation}\label{eq:strong-form}
\Lam(m) \asymp \kpi(m).
\end{equation}
\end{theorem}

\begin{conjecture}[Strong Form Conjecture]
The asymptotic equivalence~\eqref{eq:strong-form} holds under the stated regularity
conditions. The proof requires establishing a matching upper bound $\Lam(m) \lesssim
\kpi(m)$, which follows from the Lipschitz condition on $p$ and the bounded fiber
diameter via a Jensen-Shannon tensorization argument. The key steps are: the Lipschitz
condition implies that within-fiber admissibility divergence grows at most linearly with
fiber diameter; bounded fiber diameter caps the divergence; the Jensen-Shannon divergence
is bounded and symmetric and satisfies a data-processing inequality that converts
within-fiber divergence to representational mixing. This conjecture is presented as a
research target for the companion analysis.
\end{conjecture}

\section{The Regularity Conditions and Their Semantic Content}

Each of the three regularity conditions in Theorem~\ref{thm:strong} has a precise
semantic interpretation.

Bounded fiber diameter means that the ambiguity introduced by projection is bounded: the
representation does not collapse infinitely many admissibility-distinct states into one
point. This is a non-trivial requirement. It fails, for instance, when the encoding map
$\pi$ is highly lossy and many semantically incompatible states share a single
representation. In practice, bounded fiber diameter corresponds to a requirement that
the representational system has sufficient capacity to distinguish the main branches of
the admissibility structure.

Lipschitz continuity of $p(\cdot\mid x)$ in $x$ means that the admissibility structure
varies without discontinuous jumps: there are no sharp semantic phase transitions where
the accessible futures change catastrophically over a small change in position. Domains
with genuine discontinuities---sharp boundaries between legal and illegal moves in formal
games, for instance---violate this condition and require a separate analysis.

Jensen-Shannon divergence, in place of KL divergence, is required because the KL
divergence can be infinite when the support of $\pi_*\mu_U$ does not contain the support
of $\prod_i \pi_*\mu_{U_i}$. Jensen-Shannon divergence is bounded by $\log 2$,
symmetric, and metrizes weak convergence. These properties are essential for the
tensorization argument that gives the upper bound in Theorem~\ref{thm:strong}.

\section{Interpretive Consequences}

Theorem~\ref{thm:weak} establishes that collapse is a geometric necessity under
projection with hidden curvature: it is not an implementation failure, not a training
artifact, and not a property that can be eliminated by adding more representational
capacity without changing the projection geometry. Any map $\pi: X\to M$ that identifies
points in a curved admissibility manifold will produce observable mixing in proportion to
the curvature it destroys.

Theorem~\ref{thm:strong}, if the conjecture is confirmed, establishes something stronger:
collapse measurements are geometric observables. They encode curvature information up to
multiplicative constants in the same sense that gravitational lensing encodes spacetime
curvature. The mixing functional $\Lam$ is not merely bounded by $\kpi$. It is
asymptotically equivalent to it. A complete field of $\Lam$-measurements over $M$
therefore determines $\kpi$ up to a constant, and the inverse problem of recovering the
admissibility geometry from collapse data becomes well-posed in the sense of geometric
tomography.

%% ═════════════════════════════════════════════════════════════════════════════
\part{The Inverse Problem}
%% ═════════════════════════════════════════════════════════════════════════════

\chapter{Collapse as Geometric Tomography}

\section{The Inverse Problem Formulated}

The forward problem is now solved: given the admissibility manifold $(X,g)$ and the
projection $\pi: X\to M$, the Projection-Collapse Principle determines the mixing
functional $\Lam$ over $M$. The inverse problem reverses this: given a representational
manifold $M$, a field of mixing measurements $\Lam$ over $M$, and a projection $\pi$
(known or partially known), recover the projected sectional curvature $\kpi$ over the
fibers of $\pi$, and from $\kpi$, the admissibility density $p(y\mid x)$ up to the
resolution limit imposed by $\pi$.

This is an inverse problem in the sense of Tikhonov: the forward map $\kpi\mapsto\Lam$
is known from the theorem; the task is to invert it from noisy, finite-resolution
observations. Standard results from the theory of ill-posed inverse problems apply:
uniqueness of the reconstruction requires additional structural assumptions, and stability
of the reconstruction requires regularization.

\section{Analogy with Geometric Tomography}

The analogy with gravitational lensing and seismic tomography is precise enough to
import mathematical structure from those domains. In gravitational lensing, the observable
is the deflection of light passing through a gravitational field; the hidden quantity is
the spacetime metric. In seismic tomography, the observable is the travel time of elastic
waves through the earth; the hidden quantity is the elastic velocity field. In both
cases, the forward map (from the hidden geometry to the observable) is known, and the
inverse problem is to recover the geometry from measurements of the observable.

Here, the observable is the mixing functional $\Lam$ and the hidden quantity is the
projected sectional curvature $\kpi$. The forward map is the Projection-Collapse
Principle. The inverse problem has the same structure as lensing and seismology: a
geometry is inferred from its effect on the observable trajectories that pass through it.
The semantic analogue of light deflection is the distortion of representational
neighborhoods under the admissibility geometry.

\section{Resolution Limits and the Cramér-Rao Bound}

The Fisher metric on $X$ implies a fundamental resolution limit on any representation
$\pi$. The Cramér-Rao inequality, applied to the statistical manifold $(X,g)$, states
that no estimator of the position $x\in X$ from observations in $M=\pi(X)$ can achieve
variance smaller than the inverse of the Fisher information in the projected directions.
This is a geometric lower bound on the resolution of $\pi$: the representation cannot
distinguish admissibility structures more finely than the metric allows.

Below this resolution scale, the inverse problem is ill-posed: multiple distinct
admissibility geometries produce the same pattern of collapse measurements. Regularization
is required to select among the consistent solutions. The appropriate regularization
encodes a prior over admissibility geometries, and the choice of prior reflects
assumptions about the constraint structure of the domain. A maximum-entropy prior selects
the flattest geometry consistent with the data. A smoothness prior selects the most
slowly varying geometry. A constraint-consistent prior selects the geometry most
compatible with known structural properties of the domain.

\section{What Can and Cannot Be Recovered}

A precise characterization of recoverable information from $\Lam$ measurements is as
follows. The projected sectional curvature $\kpi$ is recoverable up to fiber-averaging:
the integral~\eqref{eq:kpi} can be estimated from $\Lam$ measurements, but the
point-wise curvature $|K_g(v,h)|$ at individual points in the fiber cannot be resolved
without additional measurements. The full Riemann tensor $R_g$ is not recoverable from
$\Lam$ alone: only the mixed components $K_g(v,h)$ for $v\in V_x$ and $h\in H_x$
contribute to leakage; the pure vertical and pure horizontal components require separate
probes.

The admissibility density $p(y\mid x)$ is recoverable up to the equivalence class defined
by the fiber structure of $\pi$: two density functions that agree on every fiber of $\pi$
produce the same $\Lam$ field and cannot be distinguished by collapse measurements alone.
Distinguishing within-fiber variation requires perturbation experiments---controlled
changes to the input that probe specific tangent directions---or contrastive evaluations
that compare the accessible futures of adjacent positions in the fiber.

\section{The Reconstruction Operator}

The inverse problem can be formalized as a bounded linear operator between function spaces.
Let $L^2(M, \mu_M)$ denote the space of square-integrable functions on the representational
manifold with respect to the pushforward measure $\mu_M = \pi_* \mu$. Define the
\emph{reconstruction operator}:
\begin{equation}\label{eq:recon-op}
\mathcal{R}: L^2(M, \mu_M) \to L^2(M, \mu_M),
\qquad
\mathcal{R}(\Lam) = \kpi,
\end{equation}
as the operator that maps an observed mixing field $\Lam$ to the projected sectional
curvature field $\kpi$ that produced it through the Projection-Collapse Principle.

\begin{theorem}[Stability of Reconstruction]\label{thm:reconstruction}
Under the regularity conditions of Theorem~\ref{thm:strong} (bounded fiber diameter $D$,
$L_p$-Lipschitz admissibility density, Jensen-Shannon divergence), the reconstruction
operator $\mathcal{R}$ satisfies the stability estimate:
\begin{equation}\label{eq:stability}
\|\mathcal{R}(\Lam_1) - \mathcal{R}(\Lam_2)\|_{L^2(M)}
\leq C_{D, L_p}\, \|\Lam_1 - \Lam_2\|_{L^2(M)},
\end{equation}
for a constant $C_{D,L_p}$ depending on the fiber diameter bound $D$ and the Lipschitz
constant $L_p$. Consequently, small perturbations in the observed mixing field produce
bounded perturbations in the reconstructed curvature.
\end{theorem}

\begin{proof}[Proof sketch]
From the asymptotic equivalence $\Lam(m) \asymp \kpi(m)$ of Theorem~\ref{thm:strong},
the ratio $\kpi(m)/\Lam(m)$ is bounded above and below by constants depending on $D$ and
$L_p$. Therefore:
\begin{align*}
\|\mathcal{R}(\Lam_1) - \mathcal{R}(\Lam_2)\|_{L^2}^2
&= \int_M |\kpi^{(1)}(m) - \kpi^{(2)}(m)|^2\, d\mu_M(m) \\
&\leq C^2 \int_M |\Lam_1(m) - \Lam_2(m)|^2\, d\mu_M(m)
= C^2 \|\Lam_1 - \Lam_2\|_{L^2}^2,
\end{align*}
where the inequality uses the pointwise bound $|\kpi^{(1)} - \kpi^{(2)}| \leq
C\,|\Lam_1 - \Lam_2|$ from the asymptotic equivalence. The constant $C = C_{D,L_p}$
is the one from Theorem~\ref{thm:strong}.
\end{proof}

Theorem~\ref{thm:reconstruction} establishes that the inverse problem is well-posed in
the Hadamard sense: a solution exists (by the Projection-Collapse Principle), it is
unique up to fiber equivalence (by the separating properties of $D_{\mathrm{JS}}$), and
it depends continuously on the data (by the stability estimate~\eqref{eq:stability}).
The reconstruction is therefore not merely an aspiration but a certifiably stable inverse
operation.

The stability estimate also quantifies the resolution limit of the tomographic
reconstruction: uncertainty in $\Lam$ by $\epsilon$ translates into uncertainty in $\kpi$
by at most $C_{D,L_p}\,\epsilon$. To reconstruct curvature to precision $\delta$, it
suffices to measure the mixing functional to precision $\delta / C_{D,L_p}$.

\chapter{Spectral Reconstruction Theory}

The reconstruction operator $\mathcal{R}$ of the previous chapter is shown to exist and
be stable, but its construction is left implicit. This chapter makes the construction
explicit by developing a spectral theory for the forward operator $L: \kpi \mapsto \Lam$.
The spectral decomposition transforms the reconstruction from an abstract inverse operator
into a concrete computational procedure, and identifies the ill-posedness structure that
regularization must address.

\section{The Linearized Forward Operator}

In the regime where $\Lam$ depends approximately linearly on $\kpi$---which holds in the
asymptotic regime of Theorem~\ref{thm:strong} where $\Lam \asymp \kpi$---define the
\emph{linearized forward operator}:
\begin{equation}\label{eq:forward-op}
L: L^2(M, \mu_M) \to L^2(M, \mu_M), \qquad L\kpi = \Lam.
\end{equation}
The operator $L$ encodes the geometric relationship between the hidden curvature field
and the observable mixing field. Explicitly, $L$ is an integral operator:
\begin{equation}\label{eq:forward-kernel}
(L\kpi)(m) = \int_M K(m, m')\, \kpi(m')\, d\mu_M(m'),
\end{equation}
where the kernel $K(m, m')$ describes how curvature at $m'$ affects mixing at $m$ through
the fiber geometry. The kernel is determined by the projection $\pi$ and the Fisher metric
$g$: points in the same fiber contribute directly, while points in nearby fibers contribute
through the $A$-tensor coupling~\eqref{eq:kpi-A}.

\begin{proposition}[Compactness and Self-Adjointness]\label{prop:L-props}
If the admissibility manifold $X$ is compact and the kernel $K(m,m')$ is square-integrable
over $M \times M$, then:
\begin{enumerate}
  \item[(i)] $L$ is compact as an operator on $L^2(M, \mu_M)$;
  \item[(ii)] $L$ is self-adjoint when $K(m,m') = K(m',m)$, which holds when the fiber
    geometry is symmetric (the measure $\mu_m$ is symmetric under fiber automorphisms).
\end{enumerate}
\end{proposition}

Compactness of $L$ is the precise statement of ill-posedness: the spectrum of $L$ has
zero as an accumulation point, meaning that the inverse operator $L^{-1}$ is unbounded.
Small errors in $\Lam$ can be amplified without bound by naive inversion.

\section{Singular Value Decomposition}

By the spectral theorem for compact self-adjoint operators, $L$ has a countable
orthonormal eigenbasis $\{\phi_n\}_{n \geq 1}$ of $L^2(M, \mu_M)$ with corresponding
eigenvalues $\sigma_n \geq 0$ decreasing to zero:
\begin{equation}\label{eq:svd}
L = \sum_{n=1}^\infty \sigma_n\, \phi_n \otimes \phi_n,
\qquad \sigma_1 \geq \sigma_2 \geq \cdots \geq 0, \quad \sigma_n \to 0.
\end{equation}
Expanding $\kpi = \sum_n a_n \phi_n$ and $\Lam = \sum_n b_n \phi_n$ in this basis:
\[
b_n = \sigma_n\, a_n, \qquad n \geq 1.
\]
The naive inversion $a_n = b_n / \sigma_n$ is unstable when $\sigma_n$ is small: small
measurement errors in $b_n$ produce large errors in $a_n$.

\section{Curvature Recovery via Truncated SVD}

The standard regularization strategy for compact operator inversion is truncated SVD:

\begin{definition}[Truncated SVD Reconstruction]\label{def:tsvd}
For a truncation parameter $N \geq 1$, define the truncated reconstruction operator:
\begin{equation}\label{eq:tsvd}
\mathcal{R}_N(\Lam) = \sum_{n=1}^N \frac{\langle \Lam, \phi_n\rangle_{L^2}}{\sigma_n}\, \phi_n.
\end{equation}
\end{definition}

\begin{theorem}[Reconstruction Error for Truncated SVD]\label{thm:tsvd-error}
Let $\kpi = \sum_n a_n\phi_n$ be the true curvature field and $\tilde\Lam = L\kpi +
\epsilon_\Lam$ the observed mixing field with measurement error $\|\epsilon_\Lam\|_{L^2}
\leq \epsilon$. The truncated SVD reconstruction satisfies:
\begin{equation}\label{eq:tsvd-error}
\|\mathcal{R}_N(\tilde\Lam) - \kpi\|_{L^2}^2
\leq \frac{\epsilon^2}{\sigma_N^2} + \sum_{n > N} a_n^2.
\end{equation}
The first term is the \emph{noise amplification} (decreasing in $N$) and the second is
the \emph{truncation bias} (increasing in $N$). The optimal $N$ balances these two terms.
\end{theorem}

\begin{proof}
$\mathcal{R}_N(\tilde\Lam) - \kpi = \sum_{n=1}^N \frac{\langle\epsilon_\Lam,\phi_n\rangle}
{\sigma_n}\phi_n - \sum_{n>N} a_n\phi_n$.
The squared $L^2$ norm of the first sum is bounded by $\|\epsilon_\Lam\|^2/\sigma_N^2$
(using $\sigma_n \geq \sigma_N$ for $n \leq N$); the second sum has squared norm
$\sum_{n>N}a_n^2$.
\end{proof}

\section{Tikhonov Regularization}

An alternative to truncated SVD that avoids the hard threshold is Tikhonov
regularization:

\begin{definition}[Tikhonov Reconstruction]\label{def:tikhonov}
For regularization parameter $\alpha > 0$, define:
\begin{equation}\label{eq:tikhonov}
\mathcal{R}_\alpha(\Lam) = (L^* L + \alpha I)^{-1} L^* \Lam,
\end{equation}
where $L^* = L$ (self-adjoint case). In the eigenbasis:
\[
\mathcal{R}_\alpha(\Lam) = \sum_{n=1}^\infty \frac{\sigma_n}{\sigma_n^2 + \alpha}\,
\langle\Lam, \phi_n\rangle\, \phi_n.
\]
\end{definition}

The Tikhonov filter $\sigma_n/(\sigma_n^2 + \alpha)$ smoothly suppresses small singular
values: when $\sigma_n \gg \sqrt{\alpha}$, the filter is approximately $1/\sigma_n$
(faithful inversion); when $\sigma_n \ll \sqrt{\alpha}$, the filter is approximately
$\sigma_n/\alpha$ (strong damping). The parameter $\alpha$ controls the effective
truncation.

\begin{theorem}[Tikhonov Error Bound]\label{thm:tikhonov-error}
Under the source condition $\kpi = L\omega$ for some $\omega \in L^2(M,\mu_M)$
(which holds when $\kpi$ is in the range of $L$), the Tikhonov reconstruction satisfies:
\begin{equation}\label{eq:tikhonov-error}
\|\mathcal{R}_\alpha(\tilde\Lam) - \kpi\|_{L^2}
\leq \frac{\epsilon}{2\sqrt{\alpha}} + \frac{\alpha^{1/2}}{2}\|\omega\|_{L^2},
\end{equation}
minimized at the optimal $\alpha^* = (\epsilon/\|\omega\|_{L^2})^{2/3}$, giving error
$O(\epsilon^{2/3})$.
\end{theorem}

The $O(\epsilon^{2/3})$ convergence rate is the classical Tikhonov rate for linear
inverse problems with compact operators under source conditions. It is slower than the
$O(\epsilon)$ rate for well-posed problems but is optimal for this class of ill-posed
problems.

\section{The Singular Value Spectrum as Semantic Resolution Limit}

The singular value spectrum $\{\sigma_n\}$ of $L$ encodes the resolution structure of
the reconstruction. Large singular values correspond to coarse-scale curvature features
that are reliably recovered from $\Lam$ measurements; small singular values correspond
to fine-scale features that are suppressed by the projection and difficult to reconstruct.

The rate of decay of $\sigma_n$ is determined by the geometry of the projection $\pi$
and the smoothness of the admissibility density $p(\cdot\mid x)$. For a smooth submersion
$\pi$ with compact fibers, the singular values decay polynomially: $\sigma_n \asymp
n^{-s/d}$ where $s$ is the smoothness order of $K$ and $d = \dim M$. For less regular
projections, the decay can be faster (exponential for analytic $K$) or slower.

The semantic interpretation is direct: the resolution limit of the reconstruction is
determined by the geometric structure of the projection. A projection that identifies
many admissibility-distinct states (large fibers, high $\kpi$) has a more rapidly
decaying singular value spectrum and is harder to invert. A projection that is nearly
injective (small fibers, low $\kpi$) has slowly decaying singular values and admits
accurate reconstruction. The spectral theory of $L$ therefore quantifies, in a single
compact form, the fundamental tradeoff between representational compression and
admissibility reconstruction fidelity.

\section{Triplet Measurements as Curvature Probes}

A triplet $(A, S, D)$ consists of an anchor state $A$, a valid substitute $S$, and a
misleading decoy $D$. In a faithful representation, the substitute is more similar to the
anchor than the decoy: $\mathrm{sim}(A,S) > \mathrm{sim}(A,D)$. A collapse event occurs
when this inequality reverses: $\mathrm{sim}(A,D) \geq \mathrm{sim}(A,S)$.

In the geometric framework, a triplet collapse event corresponds to a location where the
projected sectional curvature $\kpi$ is large enough that the decoy $D$ and the substitute
$S$ exchange representational proximity after projection. The curvature at the fiber over
$\pi(A)$ is sufficient to rotate the accessibility profiles of $S$ and $D$ relative to
$A$, placing $D$ closer to $A$ in $M$ than $S$ is, even though $S$ is more similar to $A$
in the admissibility geometry. Triplet evaluations are therefore point measurements of
the sign and approximate magnitude of a curvature-induced proximity distortion.

The collapse rate $\mathrm{CR} = \frac{1}{n}\sum_{i=1}^n \mathbf{1}[\mathrm{sim}(A_i,D_i)
\geq \mathrm{sim}(A_i,S_i)]$ is an estimator of the proportion of locations in $M$ where
the projected curvature exceeds the proximity threshold. A battery of triplets, distributed
across different operator classes and semantic domains, provides a coarse-grained
curvature map of the admissibility manifold.

\section{Entropy Measurements as Curvature Integrals}

The entropy of a token's neighborhood, defined as $H(t) = -\sum_{u\in N_t} q(u\mid t)
\log q(u\mid t)$ where $N_t$ is the $k$-nearest-neighbor set and $q$ is a softmax over
neighbors by cosine similarity, is an integrated curvature measurement over a
neighborhood of $t$ in $M$. High entropy indicates that the neighborhood contains points
from multiple admissibility classes---that the fiber over that region of $M$ contains
admissibility-distinct points. Low entropy indicates that the neighborhood is
admissibility-coherent.

The typed entropy $H_C(t)$, obtained by restricting the sum to neighbors of type $C$,
measures curvature within a single admissibility class. The cross-type leakage $\lambda(t)
= \sum_{C\neq C_t}\sum_{u\in N_t\cap C} q(u\mid t)$ measures the total probability mass
that has leaked across class boundaries---a direct estimate of $\Lam$ for a specific
neighborhood. The drift index $\Delta H(t) = H(t) - \mathrm{median}_{s\in R(t)} H(s)$,
where $R(t)$ is a role-matched reference set, measures the excess entropy relative to
what would be expected from the admissibility structure alone.

These quantities form a hierarchy of increasingly coarse-grained curvature observables.
Triplet measurements probe point values of proximity distortion. Entropy measurements
probe integrated curvature over neighborhoods. Surface pressure $\pi(\Omega) =
\mathbb{E}_{t\in\Omega}[\lambda(t)]$ measures average mixed curvature over an extended
region. Together, they constitute the observational toolkit for the tomographic
reconstruction of the admissibility manifold.

\section{Perturbation Experiments as Tangent Vector Probes}

Controlled perturbations of inputs---paraphrase, negation, modal substitution, tense
shift, agent-patient reversal---probe the admissibility geometry in specific tangent
directions. If a perturbation in direction $v\in T_xX$ produces large changes in the
mixing functional $\Lam$, then $v$ has large components in the vertical bundle $V_x$ and
the projection is destroying information in that direction.

More precisely, a perturbation experiment measures the directional derivative of $\Lam$
with respect to the perturbation direction: $\partial_v \Lam(\pi(x)) \approx
\frac{\Lam(\pi(x+\epsilon v)) - \Lam(\pi(x))}{\epsilon}$. Directions for which this
derivative is large are the semantically sensitive directions: small displacements in
those directions produce large changes in the admissibility structure. These are precisely
the directions of high metric weight in $g$ that the projection should preserve but
may not.

Perturbation experiments are therefore measurements of the vertical-horizontal
decomposition at specific points. They identify which tangent directions at $x$ the
representation is collapsing, and they provide gradient information that can be used to
refine the tomographic reconstruction.

\section{Combining Probes: Toward Tomographic Reconstruction}

The three classes of measurement---triplet evaluations, entropy measurements, and
perturbation experiments---provide complementary information about the admissibility
geometry. Triplets are sparse but precise: they give point estimates of proximity
distortion at specific locations. Entropy measurements are dense but coarse: they
integrate curvature over neighborhoods and cannot resolve fine-grained structure. Perturbation
experiments are directional: they probe specific tangent directions and provide gradient
information unavailable from the other two.

A tomographic reconstruction combines all three. Triplet violations locate the curvature
concentrations. Entropy gradients provide coarse-grained curvature maps. Perturbation
derivatives identify the directions in which the curvature is concentrated. Together,
under sufficient density of measurements and appropriate regularization, they constrain
the admissibility geometry up to the resolution limit imposed by the Cramér-Rao bound.

\chapter{Connection to Reentrant Value Field Theory}

\section{Bohlen's Framework and the Admissibility Manifold}

Recent work on retarded functional differential equations (RFDEs) as a foundation for
synthetic cognition provides a concrete instantiation of the abstract admissibility
geometry developed in Parts II and III~\cite{bohlen2026}. Bohlen constructs a coupled
dynamical system in which a symbolic field $H_L$ and a geometric field $X_R$, each a
section of a vertex bundle over a finite graph, are coupled through a bipartite
Hilbert-Schmidt operator with propagation delays. The central formal results---well-posedness,
compact global attractor, and delay-independent stability---are proven under explicit
admissibility conditions specified through nine synthetic design blueprints.

The connection to the present framework runs through several specific correspondences.
The history space $\mathcal{X} = C([-\tau_{\max}, 0], \mathcal{Z})$ is Bohlen's version
of the admissibility manifold $X$: it is the space of all admissible histories of the
system, equipped with the sup-norm topology. The projection from the history space to the
current state space $\pi: \mathcal{X}\to\mathcal{Z}$ given by evaluation at $t=0$ is an
instance of the general projection $\pi: X\to M$. Two history segments that agree at
$t=0$ but differ over $[-\tau_{\max}, 0)$ project to the same current state but lead to
different future dynamics---precisely the structure of the fiber $\pi^{-1}(z)$.

\section{The Small-Gain Condition as a Collapse Surface}

The stability condition $C_K^2 < \mu_L\mu_R$ in Bohlen's framework has a natural
interpretation in the admissibility geometry. Here $C_K$ is the Hilbert-Schmidt norm of
the interfield coupling operator and $\mu_L$, $\mu_R$ are the dissipativity constants of
the symbolic and geometric fields respectively. The condition is the positivity requirement
for the $2\times 2$ Schur margin matrix:
\[
\begin{pmatrix} \mu_L & -C_K \\ -C_K & \mu_R \end{pmatrix} > 0
\iff C_K^2 < \mu_L\mu_R.
\]
When this condition holds, the coupled system is globally asymptotically stable for all
finite delays. When it is violated, the coupled system loses stability.

In the language of admissibility geometry, the stability boundary $C_K^2 = \mu_L\mu_R$
is a collapse surface in the state space: a region where the admissibility geometry
becomes critically curved. Points in the history space that are near the stability
boundary have large projected sectional curvature $\kpi$: small perturbations of the
coupling strength produce large changes in the accessible futures of the system. The
small-gain condition is therefore a condition on the admissibility geometry of the
history-space projection, not merely a stability criterion in the conventional
dynamical-systems sense.

\section{The Lyapunov-Krasovskii Functional as Admissibility Density}

The Lyapunov-Krasovskii functional:
\begin{equation}\label{eq:lkf}
V(\tilde\phi) = \frac{1}{2}\|\tilde\phi_L(0)\|_F^2
  + \frac{1}{2}\|\tilde\phi_R(0)\|^2
  + \frac{C_K^2}{2\mu_L}\int_{-\tau_{R\to L}}^0 \|\tilde\phi_R(s)\|^2\,ds
  + \frac{C_K^2}{2\mu_R}\int_{-\tau_{L\to R}}^0 \|\tilde\phi_L(s)\|_F^2\,ds
\end{equation}
is a candidate for the admissibility density $\Phifield(x)$ in the neighborhood of the
equilibrium $Z^*$. It is a quadratic form on the history space that decreases along
solutions---a Lyapunov function in the classical sense. The Hessian of $V$ near the
equilibrium is well-defined and positive definite precisely because $V$ is constructed to
be convex in that region.

This is one of the locations where the Hessian construction does work, but only locally.
Near the equilibrium, where $V$ is convex, the Hessian of $V$ approximates the Fisher
metric of the admissibility density $p(y\mid x)$ restricted to the attractor. The Fisher
metric and the Hessian metric agree at the equilibrium and diverge away from it, with the
Hessian failing at saddle points and the Fisher metric remaining well-defined throughout.
Bohlen's framework provides a local approximation to the admissibility geometry near the
attractor; the Fisher metric extends this approximation globally.

\section{Blueprint Admissibility as Constraint Geometry}

The nine synthetic design blueprints in Bohlen's framework are admissibility contracts:
explicit conditions that each architectural component must satisfy for the system to
remain within the admissible region of the state space. Each blueprint specifies
regularity, boundedness, dissipativity, and symmetry conditions that, collectively,
ensure that the master RFDE satisfies the hypotheses of the formal theorems.

In the language of the present framework, the blueprint conditions are sufficient
conditions for the admissibility geometry to be non-degenerate. Blueprint 1 (symbolic
update) ensures that the symbolic field evolves within a compact, positively invariant
domain with a spectral gap---ensuring that the Fisher metric on the symbolic component of
the history space has bounded curvature. Blueprint 3 (interconnector) bounds the
Hilbert-Schmidt norm $C_K$ of the coupling kernel---bounding the mixed curvature between
the symbolic and geometric components of the history space. A violation of any blueprint
condition corresponds to a region of high $\kpi$ in the admissibility geometry: the
system leaves the region where the constraint structure supports coherent future
trajectories.

\chapter{Semantic Field Theory}

\section{From Diagnostics to Dynamics}

The inverse problem formulated in Chapter~8 recovers a static picture of the admissibility
geometry: given $\Lam$ measurements over $M$, reconstruct $\kpi$ over the fibers of
$\pi$. But semantic systems are dynamic. They evolve under optimization, learning, and
use. A complete theory requires dynamics: equations of motion for the admissibility field
under these pressures.

The central observation is that compression-driven optimization acts as a force that
flattens the admissibility field. Any optimizer that reduces representational complexity
by minimizing the distance between similar representations will, unless specifically
constrained, collapse distinctions that are not locally useful for the current task. This
is not a failure mode. It is the expected behavior of an optimizer that acts on the
representation rather than on the underlying constraint field. Collapse is what
optimization looks like from the perspective of the admissibility geometry.

\section{RSVP as a Field-Theoretic Framework}

The Relativistic Scalar-Vector Plenum (RSVP) framework provides
a candidate dynamics for the admissibility field. In RSVP, the primary objects are a
scalar field $\Phi$ and a vector field $\mathbf{v}$ over a spacetime manifold, coupled
through field equations that enforce the constraint structure of the domain. The
connection to the Fisher geometry developed in this monograph can be made precise through
the partition function construction.

\subsection*{From Admissibility Density to Gibbs Distribution}

Define the scalar admissibility potential:
\begin{equation}\label{eq:phi-log}
\Phi(x) = \log |\A(x)|,
\end{equation}
the log-volume of the admissible future cone at $x$. Given an energy function
$E(y, x) \geq 0$ measuring the cost of continuation $y$ from position $x$, define the
Gibbs-type admissibility density:
\begin{equation}\label{eq:gibbs}
p(y \mid x) = \frac{\exp(-E(y, x))}{Z(x)},
\qquad
Z(x) = \int_{\A(x)} \exp(-E(y, x))\, dy,
\end{equation}
where $Z(x)$ is the partition function. The RSVP scalar field is identified with the
log-partition function:
\begin{equation}\label{eq:phi-Z}
\Phi(x) = \log Z(x).
\end{equation}

\begin{proposition}[Fisher Metric from the Partition Function]\label{prop:fisher-Z}
For the Gibbs density~\eqref{eq:gibbs}, the Fisher information metric on the admissibility
manifold $X$ is:
\begin{equation}\label{eq:fisher-Z}
g_{ij}(x) = \partial_i \partial_j \log Z(x) - \frac{\partial_i Z(x)\, \partial_j Z(x)}{Z(x)^2}.
\end{equation}
In the case where the energy $E(y, x)$ is linear in $x$---that is, $E(y,x) = -\langle
\theta(x), \phi(y)\rangle$ for sufficient statistics $\phi(y)$ and natural parameters
$\theta(x)$---the admissibility density is an exponential family and the metric simplifies to:
\begin{equation}\label{eq:fisher-exp}
g_{ij}(x) = \partial_i \partial_j \log Z(\theta(x)),
\end{equation}
the Hessian of the log-partition function in the natural parameters.
\end{proposition}

\begin{proof}
By direct computation from the Fisher metric definition~\eqref{eq:fisher}:
\begin{align*}
g_{ij}(x)
&= \mathbb{E}_{y \sim p(\cdot\mid x)}\!\left[\partial_i \log p(y\mid x)\cdot
   \partial_j \log p(y\mid x)\right] \\
&= \mathbb{E}_{y \sim p(\cdot\mid x)}\!\left[
   \left(\partial_i(-E(y,x)) - \partial_i \log Z(x)\right)
   \left(\partial_j(-E(y,x)) - \partial_j \log Z(x)\right)\right] \\
&= \mathrm{Cov}_{y \sim p(\cdot\mid x)}\!\left[\partial_i(-E(y,x)),\, \partial_j(-E(y,x))\right],
\end{align*}
which equals the covariance of the score function, the standard information-geometric
expression~\eqref{eq:fisher-Z}. For the exponential family case $\log p(y\mid x) =
\langle\theta(x), \phi(y)\rangle - \log Z(\theta(x))$, the score function is
$\partial_i \log p = \partial_i\theta^k \cdot \phi_k - \partial_i \log Z$, and the
Fisher metric in the natural parameterization is the Hessian of $\log Z$, which is
positive definite wherever $Z$ is strictly log-convex.
\end{proof}

Proposition~\ref{prop:fisher-Z} establishes the precise bridge between RSVP field theory
and the Fisher geometry. The RSVP scalar field $\Phi = \log Z$ is not merely an analogy
for admissibility density. It is the generating function of the Fisher metric: its Hessian
in natural coordinates is the metric tensor. The RSVP field equations, which govern the
dynamics of $\Phi$, therefore govern the dynamics of the Fisher metric and, through the
Projection-Collapse Principle, the dynamics of collapse.

\subsection*{Collapse as RSVP Field Deformation}

In the RSVP framework, the field equations:
\begin{align}
\Box\Phi &= \lambda(\Phi^2 - v^2)\Phi - \nabla_\mu V^\mu, \label{eq:rsvp-scalar}\\
\nabla_\nu F^{\mu\nu} &= J^\mu[\Phi], \label{eq:rsvp-vector}
\end{align}
enforce a Mexican-hat potential for $\Phi$, producing spontaneous symmetry breaking at
locations where the admissibility density $Z(x)$ falls below the vacuum expectation $v$.
These are the semantic bottlenecks of Chapter~3: locations where the admissible future
cone contracts sharply.

Optimization that violates these equations deforms $\Phi$ away from its equilibrium
profile, flattening the partition function landscape and destroying the curvature that the
Fisher metric encodes. Compression-driven learning---gradient descent on a loss that
rewards representational compactness---is exactly such a deformation: it drives $Z(x)$
toward uniformity, reducing the variance of $p(\cdot\mid x)$ across fibers and
suppressing the mixed sectional curvature $\kpi$. The Projection-Collapse Principle
describes the observable consequence: the mixing functional $\Lam$ decreases as $\kpi$
is flattened, and the reconstruction program becomes correspondingly less informative.

This analysis suggests a dynamical interpretation of the monograph's program: the
reconstruction operator $\mathcal{R}$ of Theorem~\ref{thm:reconstruction} can be applied
not only to a static field of measurements but to a time series of measurements, tracking
the deformation of the admissibility geometry over the course of training. Collapse
curves---plots of $\Lam$ against training step---become trajectories in the space of
RSVP field configurations, and the reconstruction program recovers the underlying
field-theoretic dynamics from the observable collapse trace.



\section{Meaning as a Field-Theoretic Invariant}

The conclusion of the theoretical development can now be stated precisely. Meaning is not
a property of representations. It is an invariant of the admissibility field: a quantity
preserved under admissibility-respecting transformations and destroyed by transformations
that violate the constraint structure.

More precisely: a semantic transformation $T: X\to X$ is admissibility-respecting if it
maps admissible continuations to admissible continuations: $T(\A(x))\subseteq \A(T(x))$
for all $x\in X$. Under admissibility-respecting transformations, the Fisher metric is
preserved: $T^*g = g$. The Riemann curvature tensor is preserved: $T^*R_g = R_g$. The
Projection-Collapse Principle is preserved: the relationship between $\Lam$ and $\kpi$
holds in the transformed coordinates. Semantic integrity is the preservation of these
geometric invariants.

Semantic collapse is the observable signature of a transformation that is not
admissibility-respecting: one that maps some admissible continuations outside the
admissible region, flattening the curvature of the admissibility manifold and producing
observable mixing in the representational space. The collapse measurements $\Lam$ are the
projected residue of this geometric violation.

\chapter{Variational RSVP Geometry}

The RSVP scalar field $\Phi = \log Z$ generates the Fisher metric in the exponential
family setting (Proposition~\ref{prop:fisher-Z}). This chapter develops the full
variational formulation: the admissibility action whose Euler-Lagrange equations are the
RSVP field equations, the curvature evolution under these dynamics, and the interpretation
of symmetry breaking as collapse formation.

\section{The Admissibility Action}

Define the \emph{admissibility action functional} on the space of scalar fields
$\Phi: X \to \mathbb{R}$:
\begin{equation}\label{eq:rsvp-action}
S[\Phi] = \int_X \left(\frac{1}{2}|\nabla\Phi|_g^2 - V(\Phi)\right) d\mu_g,
\end{equation}
where $|\nabla\Phi|_g^2 = g^{ij}\partial_i\Phi\,\partial_j\Phi$ is the squared norm of
the gradient in the Fisher metric, $V: \mathbb{R} \to \mathbb{R}$ is a potential
function, and $d\mu_g$ is the Riemannian volume form of $(X,g)$.

The action~\eqref{eq:rsvp-action} has two terms. The kinetic term $\frac{1}{2}|\nabla
\Phi|_g^2$ penalizes rapid variation of the admissibility density across the manifold:
regions where $\Phi$ changes quickly contribute large kinetic energy, and the action
favors smooth admissibility profiles. The potential term $V(\Phi)$ encodes the preferred
admissibility level: the specific shape of $V$ determines the equilibrium distribution
of $\Phi$.

\section{Euler-Lagrange Equations: The Admissibility Field Equation}

\begin{theorem}[Admissibility Field Equation]\label{thm:euler-lagrange}
Critical points of the action $S[\Phi]$ satisfy the Euler-Lagrange equation:
\begin{equation}\label{eq:el}
\Delta_g \Phi = V'(\Phi),
\end{equation}
where $\Delta_g = g^{ij}\nabla_i\nabla_j$ is the Laplace-Beltrami operator on $(X,g)$.
\end{theorem}

\begin{proof}
Consider a variation $\Phi_\epsilon = \Phi + \epsilon\eta$ for compactly supported $\eta:
X \to \mathbb{R}$. Then:
\[
\frac{d}{d\epsilon}\bigg|_0 S[\Phi_\epsilon]
= \int_X \left(g^{ij}\partial_i\Phi\,\partial_j\eta - V'(\Phi)\eta\right) d\mu_g.
\]
Integrating the first term by parts (using the divergence theorem on $(X,g)$ with
compact support of $\eta$):
\[
\int_X g^{ij}\partial_i\Phi\,\partial_j\eta\, d\mu_g
= -\int_X \eta\, \Delta_g\Phi\, d\mu_g.
\]
Setting the variation to zero for all $\eta$ gives~\eqref{eq:el}.
\end{proof}

The field equation~\eqref{eq:el} is the admissibility analogue of the Laplace or
Poisson equation. It determines the equilibrium admissibility profile: at equilibrium,
the Laplacian of $\Phi$ (measuring how $\Phi$ deviates from its local average) equals
the derivative of the potential (the force pulling $\Phi$ toward the potential minimum).

\section{The Mexican-Hat Potential and Semantic Phase Transitions}

The RSVP framework uses a Mexican-hat (double-well) potential:
\begin{equation}\label{eq:mh-potential}
V(\Phi) = \frac{\lambda}{4}(\Phi^2 - v^2)^2,
\end{equation}
with $\lambda > 0$ and $v > 0$ the vacuum expectation value. This potential has two
minima at $\Phi = \pm v$ and a local maximum at $\Phi = 0$.

\begin{proposition}[Symmetry Breaking and Bottleneck Formation]\label{prop:sym-break}
For the potential~\eqref{eq:mh-potential}, the field equation~\eqref{eq:el} becomes:
\begin{equation}\label{eq:mh-el}
\Delta_g \Phi = \lambda\Phi(\Phi^2 - v^2).
\end{equation}
Regions where $|\Phi| < v$ are unstable under the dynamics: the potential force
$\lambda\Phi(\Phi^2 - v^2)$ drives $\Phi$ toward $\pm v$. Regions where $\Phi \to 0$
are semantic bottlenecks: the admissibility density collapses toward zero, and the
accessible future cone contracts to a point.
\end{proposition}

This identifies semantic phase transitions geometrically: when the admissibility field
$\Phi$ crosses the unstable region $|\Phi| < v$, the system undergoes a transition from
a semantically rich region (many accessible futures, $\Phi \approx v$) to a bottleneck
($\Phi \approx 0$). The transition corresponds to the formation of a collapse surface in
the sense of Chapter~5: a region of high projected sectional curvature that generates
large observable mixing.

\section{Fisher-RSVP Correspondence: The Metric from the Action}

The connection between the action~\eqref{eq:rsvp-action} and the Fisher metric is made
precise through the following correspondence.

\begin{theorem}[Fisher-RSVP Correspondence]\label{thm:fisher-rsvp}
For an exponential family admissibility density $p(y\mid x) = \exp(\langle\theta(x),
\phi(y)\rangle - \Phi(\theta(x)))$ with natural parameters $\theta(x)$ and log-partition
function $\Phi = \log Z$, the Fisher metric is:
\begin{equation}\label{eq:fisher-rsvp}
g_{ij}(x) = \partial_i\partial_j \Phi(\theta(x))
= \partial_i\partial_j \log Z(\theta(x)),
\end{equation}
and the Riemann curvature tensor of $g$ satisfies:
\begin{equation}\label{eq:curv-from-phi}
R_{ijkl} = \frac{1}{2}\left(\partial_i\partial_l g_{jk} + \partial_j\partial_k g_{il}
- \partial_i\partial_k g_{jl} - \partial_j\partial_l g_{ik}\right)
+ g^{pq}(\Gamma_{il,p}\Gamma_{jk,q} - \Gamma_{ik,p}\Gamma_{jl,q}),
\end{equation}
where all Christoffel symbols are computed from $\Phi$ via $g_{ij} = \partial_i\partial_j
\Phi$.
\end{theorem}

\begin{proof}
The first statement is Proposition~\ref{prop:fisher-Z}. The second statement is the
standard formula for the Riemann tensor in terms of the metric and its derivatives,
applied to $g_{ij} = \partial_i\partial_j\Phi$.
\end{proof}

Theorem~\ref{thm:fisher-rsvp} completes the RSVP-geometry bridge: the RSVP scalar field
$\Phi$ is the generating function of the Fisher metric. The action $S[\Phi]$ is therefore
an action on the space of Riemannian metrics on $X$ (within the exponential family class),
and the field equation $\Delta_g\Phi = V'(\Phi)$ is an equation of motion for the metric
tensor itself. The RSVP field theory is, in the exponential family regime, a field theory
on the space of admissibility geometries.

\section{Curvature Evolution Under RSVP Dynamics}

The field equation~\eqref{eq:el} determines how the metric $g = \partial^2\Phi$ evolves
when $\Phi$ is perturbed. Differentiating~\eqref{eq:fisher-rsvp} twice more:
\[
\frac{\partial}{\partial t} g_{ij} = \partial_i\partial_j \dot\Phi,
\]
where $\dot\Phi$ satisfies the linearized field equation $\Delta_g\dot\Phi = V''(\Phi)
\dot\Phi$. This is a Schrödinger-type equation on the admissibility manifold with
potential $V''(\Phi)$. Near the vacuum $\Phi = \pm v$, the potential is $V''(\pm v) =
2\lambda v^2 > 0$, giving a massive scalar field: small perturbations of $\Phi$ are
stable and decay exponentially. Near the bottleneck $\Phi = 0$, the potential is
$V''(0) = -\lambda v^2 < 0$, giving a tachyonic scalar field: perturbations grow
exponentially, corresponding to the onset of phase transition.

The curvature of $g$ evolves according to the linearization of~\eqref{eq:curv-from-phi}
in $\dot\Phi$. In particular, the scalar curvature $R = g^{ij}R_{ij}$ satisfies an
evolution equation that couples to the Laplacian of $V''(\Phi)$. The collapse-driving
curvature $\kpi$ is a projection of $R_{ijkl}$ and therefore inherits this dynamics.
This gives a precise sense in which the RSVP field theory predicts the temporal evolution
of collapse: as $\Phi$ evolves under~\eqref{eq:el}, the mixing functional $\Lam(t) \asymp
\kpi(t)$ evolves according to the projected curvature dynamics.

\chapter{Persistent Admissibility Topology}

The geometric theory developed in Parts II and III describes the local structure of the
admissibility manifold: curvature at points, volume of balls, geodesic distances. But
collapse and semantic integrity have global aspects that local geometry cannot capture.
Disconnected admissibility basins, topological loops that create contextual memory
effects, and voids in the accessible future space are topological features that persist
across scales. This chapter applies persistent homology to the admissibility filtration
to extract these global features.

\section{The Admissibility Filtration}

\begin{definition}[Admissibility Filtration]\label{def:filtration}
For the scalar admissibility field $\Phi: X \to \mathbb{R}$, define the
\emph{superlevel set filtration}:
\begin{equation}\label{eq:filtration}
X_\lambda = \{x \in X : \Phi(x) \geq \lambda\}, \qquad \lambda \in \mathbb{R}.
\end{equation}
As $\lambda$ decreases from $\sup\Phi$ to $-\infty$, the sets $X_\lambda$ grow
monotonically: $X_{\lambda'} \subseteq X_\lambda$ for $\lambda' \geq \lambda$.
The collection $\{X_\lambda\}_{\lambda \in \mathbb{R}}$ is the admissibility filtration.
\end{definition}

The filtration~\eqref{eq:filtration} tracks the topology of the admissibility manifold
at each threshold $\lambda$. At high $\lambda$, only the most admissible positions are
included: $X_\lambda$ consists of isolated peaks of the admissibility landscape. As
$\lambda$ decreases, basins of admissibility merge, tunnels appear, and voids form. The
homology groups $H_k(X_\lambda)$ track these topological changes.

\section{Homology of the Admissibility Filtration}

For each threshold $\lambda$, the homology groups $H_k(X_\lambda; \mathbb{Z})$ with
$k = 0, 1, 2, \ldots$ capture topological features of $X_\lambda$:
\begin{align*}
H_0(X_\lambda) &\cong \mathbb{Z}^{c(\lambda)}, \quad
  \text{where } c(\lambda) = \text{number of connected components of } X_\lambda; \\
H_1(X_\lambda) &\cong \mathbb{Z}^{\ell(\lambda)}, \quad
  \text{where } \ell(\lambda) = \text{number of independent loops in } X_\lambda; \\
H_2(X_\lambda) &\cong \mathbb{Z}^{v(\lambda)}, \quad
  \text{where } v(\lambda) = \text{number of independent voids in } X_\lambda.
\end{align*}

The semantic interpretations are:
Connected components of $X_\lambda$ are disconnected admissibility basins: regions of
semantic space with sufficient admissibility that are topologically isolated from each
other at threshold $\lambda$. A system in one basin cannot reach another by a path that
stays above the threshold.

Loops in $X_\lambda$ are cyclic admissibility structures: paths that return to the same
semantic position while remaining above threshold. These correspond to contextual memory
effects---the holonomy discussed in the open problems chapter---where returning to the
same representation along different admissibility routes yields different accessible
futures.

Voids in $X_\lambda$ are regions of admissibility space completely enclosed by the
superlevel set but interior to it: holes in the accessible future. These are the
locations where the constraint structure forbids certain continuations from all
surrounding positions.

\section{Persistence Diagrams}

\begin{definition}[Persistence Diagram]\label{def:persistence}
A topological feature (connected component, loop, or void) is \emph{born} at threshold
$\lambda_b$ and \emph{dies} at threshold $\lambda_d < \lambda_b$ (the feature appears as
$\lambda$ decreases and disappears as $\lambda$ decreases further). The \emph{persistence}
of the feature is $\lambda_b - \lambda_d$. The \emph{persistence diagram} $\mathrm{Dgm}_k$
is the multiset of birth-death pairs $(\lambda_b, \lambda_d)$ for all $k$-dimensional
features.
\end{definition}

Features with large persistence $\lambda_b - \lambda_d$ are geometrically significant:
they survive a wide range of threshold values and correspond to robust structures in the
admissibility landscape. Features with small persistence are geometrically noise: they
appear and disappear within a narrow range and correspond to local fluctuations.

\begin{theorem}[Stability of Persistence Diagrams~\cite{carlsson2009}]\label{thm:stability}
If $\Phi, \Psi: X \to \mathbb{R}$ are two admissibility fields with $\|\Phi - \Psi\|_\infty
\leq \epsilon$, then the bottleneck distance between their persistence diagrams satisfies:
\begin{equation}\label{eq:stability-pd}
d_B(\mathrm{Dgm}_k(\Phi), \mathrm{Dgm}_k(\Psi)) \leq \epsilon,
\end{equation}
for all $k$. Small perturbations of the admissibility field produce small perturbations
in the persistence diagram.
\end{theorem}

Theorem~\ref{thm:stability} means that the topological features of the admissibility
filtration are robust: they can be estimated from approximations to $\Phi$ without
catastrophic topological error, as long as the approximation error is small in the
sup-norm.

\section{Topological Bottlenecks and Disconnected Basins}

A topological bottleneck is a birth-death pair $(\lambda_b, \lambda_d)$ in
$\mathrm{Dgm}_0$ (the $H_0$ diagram): a connected component of $X_{\lambda_b}$ that
merges with another component at $\lambda_d$. The persistence $\lambda_b - \lambda_d$
measures the depth of the admissibility saddle between the two basins: to pass from one
basin to the other while staying in $X_\lambda$, one must descend at least
$\lambda_b - \lambda_d$ below the peak of the lower basin.

In semantic terms, a topological bottleneck of high persistence corresponds to two
admissibility basins that are strongly separated: semantic positions in one basin have
accessible futures that are radically different from those in the other, and moving
between them requires passing through a region of very low admissibility. This is the
global analogue of the local bottleneck of Definition~\ref{def:bottleneck}:
local bottlenecks correspond to high curvature, while topological bottlenecks correspond
to persistent $H_0$ features.

\section{Semantic Topological Invariants}

The persistence diagrams $\{\mathrm{Dgm}_k\}$ are topological invariants of the
admissibility field $\Phi$. They are preserved under admissibility-respecting
transformations (which preserve $\Phi$) and changed by transformations that deform
$\Phi$. Semantic collapse, as a deformation of the admissibility field under
compression-driven optimization, therefore changes the persistence diagrams: components
merge, loops disappear, and voids fill in as $\Phi$ is flattened.

The persistence diagrams provide a global audit of semantic integrity that complements
the local collapse diagnostics of Chapter~9. Where the collapse graph (Chapter~12)
identifies local curvature concentrations, the persistence diagram identifies global
topological structure: which admissibility basins are robustly separated, which loops
create contextual memory effects, and which voids represent globally forbidden
continuations.

The relationship between the persistence diagram and the reconstruction operator
$\mathcal{R}$ of Theorem~\ref{thm:reconstruction} is an open problem. In the geometric
tomography analogy, persistence features of the seismic velocity field can be recovered
from travel-time data; here, the question is whether $H_0$ persistence features of
$\Phi$ can be recovered from $\Lam$ measurements. Theorem~\ref{thm:stability} suggests
that they can be approximately recovered whenever $\mathcal{R}(\Lam)$ approximates
$\kpi$ closely enough to determine $\Phi$ up to the sup-norm tolerance implied by the
desired persistence resolution.

%% ═════════════════════════════════════════════════════════════════════════════
\part{Architecture and Infrastructure}
%% ═════════════════════════════════════════════════════════════════════════════

\chapter{Constraint-First Semantic Architectures}

\section{The Design Inversion}

Standard architectures for semantic systems generate representations first and apply
constraints afterward. The generation step produces a distribution over continuations,
typically by predicting next tokens from a compressed context. The constraint step filters
or reweights this distribution to enforce various properties: fluency, coherence,
factuality, safety. The constraint step is downstream of generation. It is a repair
mechanism.

The admissibility framework inverts this order. Admissibility is the primary
computational object. Generation occurs by sampling from $p(\cdot\mid x)$, the
conditional density over admissible continuations from the current semantic position
$x\in X$. The distribution $p$ encodes the constraint structure from the beginning.
There is no separate repair step because inadmissible trajectories are never entered into
the generative distribution. The constraint is the prior.

This inversion has consequences at every level of the architecture. The training objective
changes: instead of minimizing prediction error on a fixed dataset, the system learns to
approximate the admissibility density $p(y\mid x)$ of the constraint structure. The
inference procedure changes: instead of sampling from a learned next-token distribution
and filtering, the system samples from the approximated $p(\cdot\mid x)$ directly. The
evaluation protocol changes: instead of measuring accuracy against a test set, the system
is evaluated by the fidelity of its $p$-approximation, measured through triplet
evaluations, entropy diagnostics, and perturbation experiments.

\section{Admissibility Filtering as a Primitive Operation}

In a constraint-first architecture, the fundamental operation is admissibility checking:
given a proposed continuation $y$ and a current position $x\in X$, is $y\in\A(x)$? This
check is not a post-generation filter. It is the inner loop of the generative process.
All computation is organized around maintaining the current position within the admissible
region of $X$.

The difference between this and standard filtering is not superficial. A standard filter
acts on the output of a generative process that is indifferent to admissibility. It
discards inadmissible outputs after they have been generated, at a cost proportional to
the probability mass placed on inadmissible regions. A constraint-first architecture never
generates inadmissible outputs, because the generative distribution has zero support
outside $\A(x)$. The gradient signal during training is concentrated on the admissible
region. The failure modes are different: instead of generating fluent but inadmissible
content, the system risks under-representing the boundary regions of $\A(x)$ where
admissibility is uncertain.

\section{Typed Representations and Projection Discipline}

If representations are projections of admissibility geometry, they should carry
information about their position in the vertical-horizontal decomposition. A typed
representation assigns to each point $m\in M$ a label that identifies the admissibility
class $C$ to which it belongs and specifies which directions in $T_m M$ are horizontal
(reliable) and which are unreliable because they are the projections of curved fiber
directions.

Projection discipline is the requirement that inference depends only on horizontal
features: features in directions that the projection faithfully transmits from the
admissibility geometry. Inference that depends on vertical features is inference that
depends on information the representation cannot reliably encode, because the vertical
directions are subject to the curvature-induced distortions quantified by $\kpi$.
Projection discipline is therefore a structural constraint on the architecture, not a
post-hoc regularizer.

In Bohlen's framework, the analogue of projection discipline is the blueprint admissibility
system: each component of the architecture is constrained to operate within a compact,
positively invariant domain with specified dissipativity and symmetry properties. The
blueprint conditions collectively ensure that the architecture remains in the region of
the history space where the admissibility geometry is non-degenerate, and they specify
which coupling directions are reliable (those satisfying the small-gain condition) and
which are not.

\chapter{Diagnostics, Auditing, and Certification}

\section{The Diagnostic Hierarchy}

A complete diagnostic system for admissibility geometry operates at multiple scales.
Point measurements, provided by triplet evaluations, locate specific collapse events: they
identify pairs of states $(A, D)$ that have exchanged relative proximity to an anchor $A$
under the admissibility geometry. Neighborhood measurements, provided by entropy and
leakage statistics, integrate curvature over local regions: they identify areas of the
representational space where the admissibility structure is globally unstable. Regional
measurements, provided by surface pressure and collapse graphs, identify extended regions
of high mixed curvature: the collapse surfaces where multiple admissibility classes
simultaneously lose separation. Global measurements, provided by total curvature
integrals and topological invariants of the admissibility manifold, characterize the
overall geometry.

Each scale provides different information and is appropriate for different evaluation
tasks. Point measurements are appropriate for targeted audits of specific semantic
distinctions. Neighborhood measurements are appropriate for detecting general instability
in a region of the representational space. Regional measurements are appropriate for
identifying the structural weaknesses of an architecture---the locations where the
projection geometry systematically fails. Global measurements are appropriate for
comparing architectures or for tracking the evolution of admissibility geometry over the
course of training.

\section{Collapse Graphs as Audit Artifacts}

Collapse graphs---distributions of triplet violations by operator class and region---are
reinterpreted in the admissibility framework as curvature maps of the projected
admissibility geometry. A collapse graph plots, for each location in $M$ and each
admissibility class, the empirical estimate of $\kpi(m)$ derived from triplet
measurements. Regions of high collapse rate correspond to regions of high projected
sectional curvature. The class-specific structure of the graph identifies which types of
semantic distinction are most severely affected by the projection.

As audit artifacts, collapse graphs provide task-grounded, falsifiable evidence about
where a system's representational compression is destroying admissibility structure. They
are not summaries of average performance. They are maps of specific geometric failures.
A system with a high overall accuracy but a collapse graph concentrated in the modal
operator region is failing to maintain the modal distinctions on which logical inference
depends, regardless of what the accuracy metric reports.

\section{Certification via Proof Obligations}

The admissibility framework supports formal certification of constraint invariants. The
procedure mirrors the blueprint admissibility system of Bohlen's framework: export the
operator graph $G_C$ encoding the local admissibility structure, verify that it satisfies
target frame properties (reflexivity, transitivity, Euclidean closure corresponding to
the modal systems T, S4, S5), and use the verification as a certificate attached to the
model.

In practice, this can be implemented using theorem provers (Isabelle, Coq) for small
graphs exported from held-out checkpoints, or using SMT solvers (Z3) for fast screening
of frame properties encoded as first-order conditions over adjacency matrices. Counterexamples
produced by these solvers identify specific locations where the admissibility geometry
violates the target invariants. These counterexamples are hard negatives for continued
training: they identify the points in the representational space where the projection is
destroying the most curvature, and they provide gradient signal for repairing the
projection geometry.

%% ═════════════════════════════════════════════════════════════════════════════
\part{Foundations and Consequences}
%% ═════════════════════════════════════════════════════════════════════════════

\chapter{Accessibility Fields, Trajectories, and Semantic Action}

Before situating the framework within the broader literature, we develop three
complementary mathematical structures that deepen the kinematic and dynamic content of
the theory: the accessibility field as a primary object, the replacement of states by
trajectory equivalence classes, and a semantic action functional that unifies local
geometry, entropic volume, and constraint dynamics into a single variational principle.

\section{Accessibility Fields}

The admissibility manifold $X$ carries more structure than a smooth manifold with a
Riemannian metric. Associated to each point $x \in X$ is a set of admissible continuations
$\A(x) \subseteq X$. The assignment:
\begin{equation}\label{eq:access-field}
\mathcal{A}: X \to \mathcal{P}(X), \qquad x \mapsto \A(x),
\end{equation}
defines the \emph{accessibility field}: a map from semantic positions to subsets of the
manifold specifying which positions are reachable in a single admissible step. The
probability density $p(\cdot \mid x)$ over $\A(x)$ introduced in Chapter~3 is a
measurable refinement of this field, encoding not merely reachability but admissibility
weight.

The accessibility field is the primary semantic object in the framework. The local
geometry of meaning at $x$ is determined not by the position of $x$ in isolation but by
the structure of $\A(x)$: what futures are open, how densely they are accessible, and
how that structure changes as $x$ varies. Two positions $x, x' \in X$ with identical
coordinates but different accessibility fields have different meanings. Two positions with
different coordinates but identical accessibility fields are semantically equivalent.

The Fisher metric $g$ is induced by $\mathcal{A}$: it measures how rapidly the
accessibility field changes as $x$ moves. Large $g_{ij}(x)$ in direction $i$ means that
small displacements in that direction substantially reshape $\A(x)$. Metric degeneracy
means the accessibility field is insensitive to motion in that direction. The admissibility
manifold with its Fisher metric is therefore a metrized accessibility field, not merely a
Riemannian manifold equipped with a semantic interpretation.

\section{Accessibility Entropy}

The scalar admissibility field $\Phi(x) = \log Z(x)$ introduced in the RSVP chapters
encodes the log-volume of the accessible continuation space. We define the accessibility
entropy directly from $\A(x)$:
\begin{equation}\label{eq:access-entropy}
S(x) = \log \mu(\A(x)),
\end{equation}
where $\mu$ is the volume measure on $X$. This quantity measures how much future freedom
the semantic position $x$ admits: large $S(x)$ indicates many accessible continuations
and high semantic flexibility; small $S(x)$ indicates a bottleneck where the constraint
structure sharply restricts forward motion.

The accessibility entropy connects three parts of the framework. In the kinematic theory,
$S(x)$ is small precisely at the bottleneck points of Definition~\ref{def:bottleneck}:
where $V(x,r) \ll V_0(r)$, the Fisher ball is small, and the log-volume $S(x)$ is
correspondingly suppressed. In the dynamic theory, $S(x) = \Phi(x) = \log Z(x)$
identifies the accessibility entropy with the RSVP scalar field. The field equation
$\Delta_g \Phi = V'(\Phi)$ is therefore an equation governing the spatial distribution
of accessibility entropy across the manifold. In the projection theory, collapse occurs
where the projection $\pi$ maps regions of high $S$-variation into a common point in
$M$: the fiber over that point contains admissibility-distinct states with different
accessibilities, and the mixing functional $\Lam$ reflects that heterogeneity.

\section{Trajectory Primacy and Semantic Equivalence}

A semantic state $x \in X$ carries only instantaneous information: its position in the
admissibility manifold and its accessibility field $\A(x)$. Richer semantic content is
encoded in trajectories. Define a \emph{semantic trajectory} as a smooth curve $\gamma:
[0,T] \to X$ such that $\gamma(t+dt) \in \A(\gamma(t))$ for all $t$: each step of the
trajectory lies within the accessibility field of the current position.

\begin{definition}[Admissible Trajectory]\label{def:admissible-traj}
A smooth curve $\gamma: [0,T] \to X$ is an \emph{admissible trajectory} if
$p(\gamma(t+\epsilon) \mid \gamma(t)) > 0$ for all $t \in [0,T)$ and all sufficiently
small $\epsilon > 0$. The space of admissible trajectories through $x$ is:
\[
\Gamma(x) = \{\gamma: [0,T] \to X \text{ admissible} : \gamma(0) = x\}.
\]
\end{definition}

The fundamental semantic object is not the state $x$ but the \emph{trajectory equivalence
class} $[\gamma]$ of all admissible trajectories through $x$ with the same initial
accessibility profile. Two trajectories $\gamma_1, \gamma_2 \in \Gamma(x)$ are
semantically equivalent if they have the same first-order accessibility structure:
$p(\cdot \mid \gamma_1(t)) = p(\cdot \mid \gamma_2(t))$ for all $t$ in a common initial
segment.

This trajectory-first perspective resolves the fiber structure of Part~III. The fiber
$\pi^{-1}(m)$ over a representational point $m \in M$ is precisely the set of trajectory
equivalence classes that project to the same representation. Points in the fiber are not
different states; they are different trajectory histories that the projection has
identified. The semantic content of each fiber point is the trajectory class that reached
it, not the representational value assigned to it. Collapse occurs because projection
discards trajectory history, and the discarded history contains admissibility-curvature
information that the mixing functional $\Lam$ partially recovers.

\section{Constraint Basins and Semantic Objects}

A \emph{semantic object}---a concept, category, or stable meaning unit---is not a point
in $X$ or a region of $M$ but an invariant basin of the accessibility field.

\begin{definition}[Constraint Basin]\label{def:basin}
A measurable subset $B \subseteq X$ is a \emph{constraint basin} if:
\begin{equation}\label{eq:basin}
\A(x) \cap B \neq \varnothing \quad \text{for all } x \in B,
\end{equation}
i.e., every position in $B$ has at least one admissible continuation that remains in $B$.
A constraint basin is \emph{minimal} if it contains no proper invariant sub-basin.
\end{definition}

Minimal constraint basins correspond to irreducible semantic units: positions from which
the accessibility field returns the system to the basin indefinitely. The filtration
$\{X_\lambda\}$ of Chapter~16 organizes basins by admissibility level: $X_\lambda$
consists of all positions with accessibility entropy above $\lambda$, and the connected
components of $X_\lambda$ are the basins accessible at that level.

Semantic collapse at the level of concepts corresponds to basin merging: when projection
$\pi$ identifies points from two distinct basins $B_1, B_2 \subseteq X$ as a common
representation $m \in M$, the fiber $\pi^{-1}(m)$ contains trajectories from both
basins. The mixing functional $\Lam$ measures this inter-basin contamination. A
hallucination, in this framework, is a generated continuation that exits the constraint
basin of the intended concept while remaining representationally proximate to its
projection.

\section{Semantic Renormalization}

The Projection-Collapse Principle is a single-level compression theorem: it relates
the admissibility geometry of $X$ to the collapse diagnostics observable in $M$. Many
representational systems apply multiple successive compressions. The natural generalization
is a renormalization group structure on the space of admissibility manifolds.

\begin{definition}[Semantic Renormalization Group]\label{def:renorm}
A \emph{semantic renormalization} is a sequence of projections:
\begin{equation}\label{eq:renorm}
X \xrightarrow{\pi_1} M_1 \xrightarrow{\pi_2} M_2 \xrightarrow{\pi_3} \cdots
\xrightarrow{\pi_n} M_n,
\end{equation}
where each $\pi_k: M_{k-1} \to M_k$ is a smooth surjection (with $M_0 = X$) that
integrates out degrees of freedom while preserving the effective admissibility structure
relevant at scale $k$.
\end{definition}

At each level $k$, the projected sectional curvature $\kappa_{\pi_k}$ measures the
admissibility information destroyed by $\pi_k$, and the mixing functional $\Lam_k$
measures the observable collapse at that level. The Projection-Collapse Principle applies
at each level independently. The composition $\pi_n \circ \cdots \circ \pi_1$ has
projected curvature bounded by the sum of the individual curvatures (by the triangle
inequality for fiber divergence), giving a total collapse bound:
\[
\Lam_{\mathrm{total}} \leq \sum_{k=1}^n \Lam_k.
\]
This is the renormalization group analogue of the data processing inequality: each
successive compression can only increase total leakage, never decrease it. The strong
form of the Projection-Collapse Principle (Theorem~\ref{thm:strong}) applies at each
level, making the multilevel collapse computable from the single-level theory.

\section{The Semantic Action Functional}

The preceding structures --- accessibility field, accessibility entropy, admissible
trajectories, constraint basins, and renormalization --- combine into a unified dynamical
principle through the semantic action functional.

\begin{definition}[Semantic Action]\label{def:action}
For an admissible trajectory $\gamma \in \Gamma(x_0)$ of duration $T$, define the
\emph{semantic action}:
\begin{equation}\label{eq:sem-action}
\mathcal{S}[\gamma] = \int_0^T \mathcal{L}(\gamma(t), \dot\gamma(t))\, dt,
\end{equation}
where the \emph{semantic Lagrangian} is:
\begin{equation}\label{eq:lagrangian}
\mathcal{L}(x, v) = \frac{1}{2} g_x(v, v) - S(x)
= \frac{1}{2} g_{ij}(x)\, v^i v^j - \log\mu(\A(x)).
\end{equation}
\end{definition}

The Lagrangian~\eqref{eq:lagrangian} has two terms. The kinetic term $\frac{1}{2}g_x(v,v)$
is the squared Fisher speed of the trajectory: it penalizes rapid deformation of the
admissibility profile. The potential term $-S(x) = -\log\mu(\A(x))$ is the negative
accessibility entropy: it penalizes positions of low admissibility volume, attracting
trajectories toward regions of high semantic flexibility. Trajectories minimizing
$\mathcal{S}[\gamma]$ subject to fixed endpoints are the solutions to the
Euler-Lagrange equation:
\begin{equation}\label{eq:sem-el}
\ddot{x}^k + \Gamma^k_{ij}\dot{x}^i\dot{x}^j = \frac{1}{2}g^{kl}\partial_l S(x),
\end{equation}
which is the geodesic equation~\eqref{eq:geodesic} with an additional force term driving
trajectories toward regions of high accessibility entropy. In the limit $S(x) = \mathrm{const}$,
equation~\eqref{eq:sem-el} reduces to the pure geodesic equation~\eqref{eq:geodesic}.

The semantic action connects all components of the framework. The Fisher metric supplies
the kinetic term and the local geometry. The accessibility entropy supplies the potential
and identifies bottlenecks as regions of suppressed volume. Constraint basins are the
attractors of the dynamics: trajectories minimizing $\mathcal{S}$ are drawn toward the
high-entropy cores of each basin. The projection collapse diagnostics measure the extent
to which the projected representation of an action-minimizing trajectory has lost its
underlying constraint structure. And the RSVP field equation $\Delta_g\Phi = V'(\Phi)$
emerges as the static limit of the action-minimizing condition when $S = \Phi$ and the
trajectory is a constant-speed geodesic: the spatial distribution of accessibility entropy
satisfies a field equation that balances diffusion against the potential force.

\chapter{Intellectual Foundations and Related Frameworks}

The framework developed in this monograph sits at the intersection of several established
mathematical and scientific traditions. This chapter positions the theory relative to
those traditions, both to acknowledge the work it builds on and to make precise what is
new: the geometric theory of admissibility is not a semantic embedding model, a modal
logic, a latent-space model, an information bottleneck, or an optimal transport theory.
It is a constraint-first field theory from which partial aspects of each of those
frameworks emerge as special cases or limiting regimes.

\section{Information Geometry}

The Fisher information metric is the mathematical foundation of the admissibility geometry.
The systematic theory of statistical manifolds equipped with the Fisher metric is due to
Amari~\cite{amari2016information,amari2000methods}, building on the invariance theorem of
Chentsov~\cite{cencov1982}. The $\alpha$-connections, dual geometries, and the role of
the exponential and mixture families in information geometry are developed in full
generality in~\cite{amari2000methods}.

The present framework uses information geometry in a non-standard way. In standard
applications, the statistical manifold parameterizes a family of probability distributions
over observed data, and the Fisher metric measures distinguishability of models. Here,
the manifold $X$ parameterizes semantic positions, and the distributions $p(\cdot \mid
x)$ are over admissible continuations rather than data. The Fisher metric measures
semantic distinguishability: how different the accessible futures are at neighboring
positions. The invariance theorem (Theorem~\ref{thm:fisher-unique}) is imported directly
from Chentsov's result, re-interpreted in the semantic setting.

\section{Riemannian Geometry and Submersion Theory}

The differential geometry of Riemannian manifolds and smooth maps provides the language
for the projection analysis of Part~III. Standard references include Lee~\cite{lee2018}
for an introduction to Riemannian manifolds and O'Neill~\cite{oneill1983} for semi-Riemannian
geometry and the theory of submersions. The O'Neill tensors $A$ and $T$, the curvature
decomposition formulas~\eqref{eq:app-HH}--\eqref{eq:app-VV}, and the relationship between
total space, base, and fiber curvature are developed in O'Neill's original
paper~\cite{oneill1966} and the monograph~\cite{oneill1983}.

The Bishop-Gromov comparison theorem (Theorem~\ref{thm:bishop-gromov}) is classical; see
for instance Cheeger and Ebin~\cite{cheeger1975} or Petersen~\cite{petersen2016} for
modern treatments. The application to accessibility volume growth in Chapter~6 is new.

\section{Optimal Transport}

The Wasserstein geometry of probability measures is deeply connected to the accessibility
structure of the admissibility manifold. Villani~\cite{villani2009,villani2003} develops
the optimal transport framework in full generality. The Wasserstein-1 distance appears
in the discrete curvature formula~\eqref{eq:ollivier} for Ollivier-Ricci curvature,
which is developed in~\cite{ollivier2009}.

The accessibility entropy and the geodesic distance $d_F$ have natural optimal transport
interpretations: $d_F(x_1, x_2)$ measures the minimum cost of transporting the
admissibility distribution at $x_1$ to that at $x_2$, in the Fisher-Rao sense. The
curvature-divergence inequality~\eqref{eq:js-lip} can be read as a transport cost
estimate. The renormalization group structure of Section~17.5 is related to the
coarse-graining of transport plans under successive marginalizations.

The present framework is not an optimal transport theory: it does not compute transport
plans or solve Monge-Kantorovich problems. But optimal transport provides a natural
language for several of its components, and the Wasserstein geometry of $\mathcal{P}
(\Omega)$ is the ambient space in which the admissibility manifold is immersed via
$\Psi: X \to \mathcal{P}(\Omega)$ (Proposition~\ref{prop:immersion}).

\section{Information Bottlenecks and Compression}

The information bottleneck method of Tishby, Pereira, and Bialek~\cite{tishby2000}
seeks a compressed representation of an input variable $X$ that preserves maximal
information about a target variable $Y$. The deep variational information bottleneck
of Alemi et al.~\cite{alemi2017} extends this to deep networks via variational inference.

The Projection-Collapse Principle can be interpreted as a geometric generalization of
bottleneck theory. In the bottleneck framework, the quality of a compression is measured
by the mutual information it preserves. In the present framework, quality is measured by
the projected sectional curvature $\kpi$ it destroys: a compression that preserves all
mixed curvature is semantically faithful; one that destroys mixed curvature creates
observable mixing. The two frameworks agree in spirit --- both treat compression as
necessarily lossy and ask which losses matter --- but differ in their mathematical object:
mutual information is a scalar; projected curvature is a field.

\section{Topological Data Analysis}

Persistent homology and its applications to data analysis are developed in Carlsson~\cite{carlsson2009}
and the survey of Edelsbrunner and Harer~\cite{edelsbrunner2010}. The admissibility
filtration of Chapter~16 and the stability theorem for persistence diagrams
(Theorem~\ref{thm:stability}) are direct applications of this machinery to the
admissibility scalar field $\Phi$.

The connection to sheaf theory, developed briefly in the open problems chapter, is treated
in Bredon~\cite{bredon1997} and Robinson~\cite{robinson2014}. The cohomological
obstruction $H^1(X, \mathcal{A})$ for a sheaf of admissibility structures to glue
globally is the formal language for the local-to-global consistency conditions that
appear informally throughout the monograph.

\section{Geometric Deep Learning}

Bronstein et al.~\cite{bronstein2021} survey the geometric deep learning framework,
which seeks to build neural architectures that respect the symmetries and geometry of
their input domains. The reentrant value field architecture of Bohlen~\cite{bohlen2026},
discussed in Chapter~13, is an instance of this program applied to coupled
symbolic-geometric systems.

The present framework is related but orthogonal: rather than building architectures that
respect a given geometry, it studies the geometry that arises from the constraint
structure of a semantic domain and asks what any architecture must preserve to be
semantically faithful. Geometric deep learning specifies the geometry from outside;
admissibility geometry derives the geometry from inside.

\section{Constraint-Based Accounts of Cognition}

Several traditions in cognitive science and theoretical neuroscience treat cognition as
navigation through constrained state spaces. The free-energy principle of Friston~\cite{friston2010}
models perception and action as joint minimization of a variational free energy that
bounds surprise. The predictive processing framework of Clark~\cite{clark2016} interprets
cognition as hierarchical prediction error minimization under generative models.

The admissibility framework is related to these approaches in treating cognition as
constraint-governed, but differs in the direction of ontological priority. Free-energy
and predictive processing accounts take the generative model as given and derive behavior
from it. The admissibility framework takes the constraint structure as primary and derives
the generative model from it as the density $p(y \mid x)$ over admissible continuations.
The two approaches make different predictions about what happens when representations
fail: in free-energy accounts, failure is prediction error; in the admissibility
framework, failure is exit from the admissible region, measured geometrically by
projected curvature. These are not equivalent descriptions.

\section{Variational Field Theory}

The admissibility action~\eqref{eq:sem-action} and the Euler-Lagrange equations of
Chapters~15 and~17 are instances of variational field theory on a Riemannian manifold.
Standard treatments include Marsden and Ratiu~\cite{marsden1999} for the mechanics of
constrained systems on manifolds and Bleecker~\cite{bleecker1981} for gauge theory and
variational principles.

The RSVP field equations~\eqref{eq:rsvp-scalar}--\eqref{eq:rsvp-vector} are analogous
to the coupled scalar-vector field theories of mathematical physics, with the Mexican-hat
potential~\eqref{eq:mh-potential} playing the role of the Higgs potential. The
spontaneous symmetry breaking and vacuum structure of that potential, applied to the
admissibility field, produce the bottleneck formation and semantic phase transitions of
Chapter~15.

\section{Process Philosophy and Structural Ontology}

The trajectory-first and constraint-first orientation of the framework connects to a
long tradition in process philosophy. Whitehead's process and reality~\cite{whitehead1929}
treats events and processes as ontologically prior to substances and states; the present
framework treats transitions and admissibility as prior to positions and representations.
Bergson's creative evolution~\cite{bergson1911} emphasizes the primacy of duration and
continuous becoming over static spatial representation; the admissibility manifold
formalizes this intuition as a Riemannian geometry on the space of constraint-governed
transitions.

These connections are philosophical rather than technical. The admissibility framework
does not import the metaphysics of process philosophy; it shares the structural commitment
that the primary objects of a semantic theory should be transitions rather than states,
and that meaning should be grounded in the lawfulness of continuation rather than the
fidelity of representation.

\section{Positioning Relative to Existing Frameworks}

The admissibility geometry is not any of the following: a semantic embedding model
(which represents meaning by position in a vector space and measures semantic similarity
by geometric proximity); a modal logic (which models necessity and possibility through
accessibility relations between possible worlds, without a metric or curvature structure);
a knowledge graph (which represents semantic relations as labeled edges in a discrete
network, without a differentiable geometry); a latent-space model (which compresses
observations into a low-dimensional representation space, without reference to constraint
structure); a free-energy theory (which grounds cognition in the minimization of
variational free energy under a generative model); an optimal transport theory (which
computes plans for redistributing probability mass at minimum cost); or an information
bottleneck framework (which optimizes the tradeoff between compression and task-relevant
information).

From each of these, partial aspects can be recovered as special cases or limiting regimes.
The Fisher metric reduces to the Euclidean metric in flat regions of the admissibility
manifold, recovering the proximity-as-similarity intuition of embedding models. The
accessibility field on a discrete graph with Boolean admissibility reduces to a Kripke
frame, recovering possible-world accessibility. The renormalization group hierarchy
reduces to successive information bottleneck compressions when the admissibility structure
is uniform across levels. The semantic action reduces to free-energy minimization in the
limit where the accessibility entropy is identified with negative variational free energy.

What the admissibility geometry adds to each of these is the geometric layer: the
curvature of the constraint field, the fiber structure of the projection, the mixing
functional as a geometric observable, and the reconstruction program that recovers hidden
curvature from collapse data. These are not present in any of the existing frameworks,
and it is their combination that constitutes the specific contribution of the present work.



\section{The Ontological Status of Admissibility}

Before developing the philosophical consequences of the framework, we address a
foundational question that any rigorous reader will pose: what kind of object is the
admissibility distribution $p(y \mid x)$? Is it a learned probability model, an empirical
frequency, a constraint density, or something more fundamental?

The answer matters because the entire geometric construction rests on $p$. If $p$ is
simply another representation, the theory may be circular. If $p$ is something
structurally different from representations, the theory escapes circularity.

Representations and admissibility structures occupy fundamentally different ontological
levels. A representation describes states: it assigns labels, coordinates, or vectors to
positions in a domain. An admissibility structure describes transformations between states:
it specifies which continuations remain lawful from each position. The distinction is not
terminological. Formally, let $r: \Omega \to M$ be a representation map assigning labels
to states. The admissibility structure is characterized by the family $\{\A(x)\}_{x \in
X}$ specifying which continuations are structurally valid from $x$. While multiple
representations may encode the same admissibility structure, the converse does not hold
in general: different admissibility structures produce geometrically distinct manifolds
$(X, g)$ that no relabeling of states can identify.

This asymmetry is the key. The admissibility field is an invariant under representational
change: if $r$ is replaced by any other representation $r'$ that encodes the same
continuation constraints, the admissibility geometry is preserved. Representations are
coordinate systems imposed upon a deeper field of lawful transitions. The admissibility
density $p(y \mid x)$ is not a learned distribution that happens to be useful; it is the
encoding of which transitions are structurally valid and with what constraint force.
Systems that infer $p$ from data are approximating this structural field from evidence,
in the same sense that physical measurements approximate the metric tensor of spacetime
from geodesic observations. The approximation may be imperfect; the target it approximates
is not itself another representation.

This establishes the ontological ground of the theory. Meaning is not grounded in labels
but in the preservation of lawful continuation. The admissibility field generates
geometry; representations are secondary projections of that geometry; collapse is the
observable signature of a projection that destroys curvature the geometry contains.

\section{Kinematics and Dynamics of Admissibility}

The framework developed in this monograph naturally separates into two complementary
layers, and making this separation explicit clarifies the relationship between the
geometric and field-theoretic components of the theory.

\begin{definition}[Semantic Kinematics]
\emph{Semantic kinematics} is the study of the geometric structure of the admissibility
manifold $(X, g)$ independent of its temporal evolution: the curvature, geodesics, volume
growth, fiber structure, and projection geometry that characterize admissibility at a
given moment.
\end{definition}

\begin{definition}[Semantic Dynamics]
\emph{Semantic dynamics} is the study of how admissibility structures evolve under the
RSVP field equations, compression-driven optimization, and constraint flows: the temporal
development of $\Phi$, the formation and dissolution of bottlenecks, and the evolution
of collapse surfaces.
\end{definition}

Parts II and III of this monograph are primarily kinematic. They answer the question:
given an admissibility structure, what is its geometry, and what does projection do to
that geometry? The RSVP chapters of Part IV are primarily dynamic. They answer the
question: how does the admissibility field evolve, and what is the field-theoretic law
governing that evolution?

The relationship between these layers mirrors the relationship between geometry and
dynamics in physics. The metric tensor of spacetime describes its geometric structure
--- the curvature of the arena in which trajectories exist. The Einstein field equations
describe how matter and energy cause that geometry to evolve. Neither layer is
reducible to the other: kinematics tells us the shape of possibility; dynamics tells
us how that shape changes.

In the present framework: admissibility geometry provides the static architecture of
semantic possibility, encoding which transitions are lawful and at what curvature.
RSVP dynamics supplies the equations governing how that architecture evolves under
compression, optimization, and learning. The Projection-Collapse Principle is a kinematic
theorem: it holds for any admissibility geometry at any moment. The prediction of how
collapse evolves during training requires the dynamic layer: the RSVP field equations
determine the trajectory of $\Phi(t)$, and through Theorem~\ref{thm:fisher-rsvp},
the trajectory of the Fisher metric $g(t)$ and the collapse functional $\Lam(t)$.

\section{Semantic Event Horizons}

The reconstruction program of Part IV recovers the projected curvature $\kpi$ from
collapse measurements $\Lam$. This recovery has a fundamental limit that deserves
explicit treatment as a structural feature of the theory rather than a technical
inconvenience.

The Jensen-Shannon divergence satisfies $0 \leq D_{\mathrm{JS}}(p, q) \leq \log 2$ for
all distributions $p$ and $q$. The mixing functional $\Lam$, defined using
$D_{\mathrm{JS}}$, is therefore bounded above by $\log 2$ regardless of the magnitude of
the underlying curvature. As $\kpi(m)$ grows without bound, $\Lam(m)$ eventually
saturates near its maximum value and ceases to carry curvature information:

\begin{equation}\label{eq:event-horizon}
\lim_{\kpi(m) \to \infty} \frac{\partial \Lam(m)}{\partial \kpi(m)} = 0.
\end{equation}

Beyond the threshold where this saturation sets in, distinct curvature values produce
observationally indistinguishable mixing signals. No finite collection of measurements
of $\Lam$ can distinguish a region of curvature $\kpi = C$ from one of curvature
$\kpi = 2C$ once both exceed the saturation threshold.

\begin{definition}[Semantic Event Horizon]\label{def:event-horizon}
A \emph{semantic event horizon} at $m \in M$ is a regime in which $\partial\Lam(m)/
\partial\kpi(m) \approx 0$: the mixing functional is saturated and further increases in
hidden curvature are undetectable from collapse measurements alone.
\end{definition}

The existence of semantic event horizons is not a defect of the framework. It is a
structural theorem about the limits of any observational program that measures
admissibility geometry from the representational side. The bound $\Lam \leq \log 2$ is
tight: there exist configurations where $\kpi$ is large but $\Lam$ is at its maximum,
and no measurement of $\Lam$ alone can determine how far beyond the threshold the
curvature has grown.

This establishes a principled ceiling on auditability: geometric features of the
admissibility manifold that lie beyond the semantic event horizon cannot be reconstructed
from collapse data. They require either direct access to the admissibility field (which
may be unavailable), supplementary measurements from a different observable (such as
perturbation experiments probing specific tangent directions), or structural assumptions
that constrain the geometry beyond what $\Lam$ alone can determine.

The horizon is movable. Using alternative divergences with unbounded range---Rényi
divergences of order $\alpha > 1$, for instance---the saturation threshold can be pushed
to larger values of $\kpi$, extending the reconstruction regime at the cost of less
tractable analysis. The bounded Jensen-Shannon divergence is the natural choice for the
core theory because of its metric properties; the extension to unbounded divergences is
an open technical program.

\section{Discrete Admissibility Geometry}

The smooth differential geometry developed throughout this monograph applies to semantic
systems with continuous parameterizations: probability families that vary smoothly with
position, manifolds with well-defined tangent spaces, and projections that are smooth
submersions. Many domains of practical significance are not continuous in this sense.
Programming languages, formal proof systems, legal codes, and symbolic reasoning
architectures have transition structures that are combinatorial: admissibility is
determined by syntactic rules, type constraints, or logical entailment, not by a smooth
density function.

The continuum framework applies to such systems only as an approximation, valid when
the density of admissible transitions is high enough that discrete structure can be
idealized as smooth. Where this approximation fails, a discrete version of admissibility
geometry is required.

Let $G = (V, E)$ be a directed admissibility graph. Vertices $v \in V$ represent
semantic states and directed edges $(v, w) \in E$ represent admissible transitions.
The accessibility structure is encoded by the neighborhood:
\[
\A(v) = \{w \in V : (v, w) \in E\},
\]
with a probability measure $p(\cdot \mid v)$ supported on $\A(v)$ that assigns weights
to admissible transitions. The admissibility manifold becomes a weighted directed graph,
and the Fisher metric becomes the matrix:
\[
g_{ij}(v) = \sum_{w \in \A(v)} \partial_i \log p(w \mid v) \cdot \partial_j \log p(w
\mid v) \cdot p(w \mid v),
\]
where the derivatives are with respect to parameters of the transition distribution.

Curvature in this discrete setting is captured by Ollivier-Ricci curvature~\cite{ollivier2009},
which measures the contraction or expansion of Wasserstein distance under the random
walk defined by $p$. For two adjacent vertices $v_1, v_2 \in V$:
\begin{equation}\label{eq:ollivier}
\kappa_{\mathrm{OR}}(v_1, v_2) = 1 - \frac{W_1(p(\cdot\mid v_1),\, p(\cdot\mid v_2))}{d(v_1, v_2)},
\end{equation}
where $W_1$ is the Wasserstein-1 distance and $d$ is the graph distance. Positive
Ollivier-Ricci curvature indicates that the neighborhoods of $v_1$ and $v_2$ overlap
substantially (the random walks from adjacent vertices spread toward each other);
negative curvature indicates that neighborhoods diverge. This is the discrete analogue
of the sectional curvature discussion in Chapter~6.

The projection $\pi: X \to M$ becomes a graph homomorphism or, in the sheaf-theoretic
generalization, a morphism of sheaves over the admissibility graph. The Projection-Collapse
Principle applies in this discrete setting with $\kpi$ replaced by Ollivier-Ricci
curvature and $\Lam$ replaced by the total variation or Wasserstein mixing between
pushforward distributions on the quotient graph.

The smooth and discrete theories share the same conceptual architecture: admissibility
structure is primary, representations are projections, curvature measures the information
destroyed by projection, and collapse diagnostics recover that curvature from the
observable side. The continuum theory emerges as a limit of the discrete theory when
the admissibility graph becomes dense and the transition distributions vary smoothly.
This limit is analogous to the relationship between discrete lattice field theories and
their continuum limits in physics: the discrete version is the more fundamental object,
and the smooth manifold is an effective description valid at sufficiently large scales.

\section{Meaning Without Representation}

The central philosophical claim of this monograph is that meaning is not a property of
representations. It is an invariant of the admissibility field---a quantity preserved
under admissibility-respecting transformations and destroyed by transformations that
violate the constraint structure. Representations are projections of that field. They are
real, causally efficacious, and practically indispensable. But they are not the primary
objects. The field is prior.

This is not a form of anti-representationalism in the philosophical sense. The claim is
not that representations do not exist or that they play no role in cognition. The claim
is ontological: constraint is prior to representation. The admissibility field generates
the geometry from which representations are derived. A representation that faithfully
projects the admissibility geometry preserves meaning. One that destroys curvature loses
it. The degree of preservation is measured by the Projection-Collapse Principle.

The practical consequence is a design principle: semantic systems should be built
constraint-first, not representation-first. The admissibility structure should be
specified, approximated, or learned before the representational compression is applied.
Adding constraints after the fact, as a downstream repair to a representation that was
built without them, is architecturally analogous to adding epicycles: it corrects the
specific failure modes that motivated the addition but does not address the underlying
geometric cause.

\section{The Chinese Room Revisited}

Searle's Chinese Room argument holds that formal symbol manipulation does not constitute
understanding because it lacks intentionality: the symbols are manipulated according to
rules but they do not mean anything to the system doing the manipulation~\cite{searle1980}.
The standard response to this argument within AI research is to add more formal structure:
more sophisticated rules, richer representations, more comprehensive coverage of the
relevant distinctions.

The admissibility framework suggests a different response. The Chinese Room fails not
because it lacks the right symbols but because it operates on representations rather than
on the admissibility field those representations project from. The rulebook that the Room
operator consults is a description of admissibility relations, not the admissibility
relations themselves. A Modal Proofing Kernel adds more pages to the rulebook. It does not
move the operator from the representational manifold $M$ back to the admissibility
manifold $X$.

Searle would likely agree with this diagnosis, though not with the remedy proposed here.
His own view is that intentionality requires biological substrate. The admissibility
framework is agnostic on this question: it does not claim that recovering the
admissibility geometry from a compressed representation constitutes understanding in the
full philosophical sense. It claims only that the geometric recovery is a necessary
condition for the kind of semantic fidelity that logical inference requires. Whether it is
sufficient is a deeper question that the present framework leaves open.

\section{Indexicality and the Admissibility Manifold}

Indexical expressions---``I,'' ``here,'' ``now''---are not reducible to descriptions
because their semantic content depends on the position of the speaker within the
constraint field, not on any property of the token
type~\cite{perry1993,kaplan1989}. The same token ``I'' expresses different content when
uttered from different positions in the admissibility manifold, because the admissible
continuations available to different agents differ.

The admissibility manifold provides a natural home for indexical content. The deictic
anchor of an indexical is a point $x\in X$: the current position of the utterance within
the constraint field. The content of the indexical is a function on $X$ that varies with
$x$: it picks out the agent, location, or time whose position in $X$ is the anchor. The
collapse of indexical meaning in representational systems is the projection of this
position-dependent function onto a position-independent representation. Any representation
that does not encode the anchor point $x$ as part of the semantic state will flatten
indexicals into descriptions, producing the conflation of ``I'' with ``the speaker'' that
is characteristic of indexical collapse.

The fix is not to add a special mechanism for indexical handling. The fix is to include
the current position in $X$ as part of the representational state: to ensure that the
projection $\pi$ is equivariant with respect to the group of admissibility-preserving
transformations, so that the anchor point is preserved under changes of context.

\chapter{Open Problems}

\section{Foundational Problems}

The most important open problem is the proof of the Strong Form of the
Projection-Collapse Principle. The conjecture states that $\Lam(m)\asymp\kpi(m)$ under
bounded fiber diameter, Lipschitz continuity of $p(\cdot\mid x)$, and Jensen-Shannon
divergence. The weak form is established by the argument of Chapter~6. The strong form
requires a matching upper bound: $\Lam(m)\lesssim\kpi(m)$. The proof strategy involves
a Jensen-Shannon tensorization argument combined with the Lipschitz bound on $p$ and the
bounded fiber diameter, but the details require careful treatment of the measure on the
fiber and the relationship between within-fiber divergence and representational mixing.

A second foundational problem is the characterization of the admissibility manifolds that
arise from natural language, formal reasoning, and perceptual systems. The general
construction of $X$ from $p(y\mid x)$ is domain-independent. For specific domains, the
geometry of $X$ reflects the specific constraint structure of the domain. Characterizing
this geometry---identifying the curvature concentrations, topological invariants, and
symmetry groups---is a program of empirical differential geometry applied to semantic
systems.

\section{Geometric Problems}

The O'Neill tensor for the submersion $\pi: X\to M$ encodes the mixed curvature in a
form amenable to explicit computation~\cite{oneill1966}. The O'Neill formulas express the
sectional curvature of the base $M$, the fiber $\pi^{-1}(m)$, and the mixed sectional
curvature $K_g(v,h)$ in terms of the A-tensor and T-tensor of the submersion. Developing
these formulas for the specific case where $g$ is the Fisher metric and $\pi$ is a neural
encoder would give a computationally tractable route to estimating $\kpi$ without relying
on triplet measurements.

A related problem is the characterization of the admissibility manifolds arising from
specific constraint structures. Boolean constraint satisfaction produces a polytope;
linear constraint satisfaction produces a polytope with a natural metric structure from
the Fisher information of the feasibility distribution. More complex constraint
structures---temporal, modal, causal---produce geometries that have not been
systematically characterized.

\section{Dynamical Problems}

The dynamics of the admissibility field under compression-driven optimization is a central
open problem. The forward direction is clear: optimization that acts on the representation
$M$ rather than the admissibility manifold $X$ flattens the admissibility geometry by
collapsing distinctions that are not locally useful. But the precise equations of motion
for the admissibility field under gradient descent, contrastive learning, or reinforcement
learning have not been derived.

The RSVP framework provides a candidate dynamics, but the connection between the RSVP
field equations and the gradient flow on the Fisher metric remains to be made explicit.
The question is whether there exists an action functional whose Euler-Lagrange equations
are both the RSVP field equations and the natural gradient flow on the space of
admissibility densities. If such an action exists, it would unify the constraint-field and
optimization-theoretic approaches in a single variational framework.

\section{Computational Problems}

Efficient algorithms for estimating $\kpi$ from finite measurements are needed for
practical application of the framework. The current approach---triplet evaluations,
entropy measurements, perturbation experiments---provides qualitative evidence about the
curvature structure but does not produce quantitative estimates suitable for the
tomographic reconstruction of $p(y\mid x)$.

A more systematic approach would apply the tools of topological data analysis, in
particular persistent homology, to the collapse graph data. The barcode of a filtration
of the collapse graph by collapse rate would provide a multi-scale summary of the
curvature structure, identifying topological features of the admissibility manifold at
different resolution scales. This approach has been developed for data manifolds in
general~\cite{carlsson2009}, but its application to admissibility manifolds specifically
requires adapting the filtration to the semantics of the domain.

\chapter*{Conclusion: Structure Must Precede Surface}
\addcontentsline{toc}{chapter}{Conclusion}

The argument of this monograph can be stated without the mathematical apparatus in a
single paragraph. Meaning is the structure of admissible transitions through a hidden
constraint field. Representations are projections of that structure. Collapse is what
happens when projection destroys curvature the field contains. Collapse diagnostics are
measurements of that curvature from the projected side. The inverse problem---recovering
the hidden geometry from its observable trace---is the central task of a semantics-aware
account of linguistic and cognitive systems.

The mathematical apparatus establishes that this is not a metaphor. The Fisher information
metric on the admissibility manifold is positive semi-definite everywhere, well-defined at
the locations where Hessian-based approaches fail, and encodes a notion of semantic
distinguishability grounded in the structure of admissible transitions. The
vertical-horizontal decomposition of the projection identifies the mixed sectional
curvature $\kpi$ as the relevant geometric quantity: the curvature that projection
destroys and that collapse diagnostics partially recover. The Projection-Collapse Principle
establishes the inequality $\Lam\geq f(\kpi)$ as a theorem and the equivalence
$\Lam\asymp\kpi$ as a well-posed conjecture with a clear proof strategy.

The broader lesson is not specific to language models, embedding spaces, or modal logic.
It is a lesson about the relationship between constraint, geometry, and representation in
any system that encodes structured information in compressed form. Observable failures of
representation contain information about hidden constraint structure. That information is
recoverable, at least partially, through the systematic analysis of collapse. The
geometry can be reconstructed from the trace it leaves when it is destroyed.

This is not a theory about what is missing from current systems. It is a theory about
what is present in the collapse patterns those systems already exhibit. The collapse is
the message.

%% ═════════════════════════════════════════════════════════════════════════════
\appendix
%% ═════════════════════════════════════════════════════════════════════════════

\chapter{Information Geometry: Background and Notation}

The Fisher information metric arises in a broader context developed systematically by
Amari and colleagues~\cite{amari2016information}. A statistical manifold is a smooth
manifold $\mathcal{S}$ whose points parameterize a family of probability distributions
$\{p(\cdot;\theta): \theta\in\mathcal{S}\}$. The Fisher information metric is the
Riemannian metric on $\mathcal{S}$ given by:
\[
g_{ij}(\theta) = \mathbb{E}_{x\sim p(\cdot;\theta)}\!\left[
  \frac{\partial\log p(x;\theta)}{\partial\theta^i}
  \frac{\partial\log p(x;\theta)}{\partial\theta^j}
\right].
\]
In the admissibility manifold of this monograph, the parameter $\theta$ is the semantic
position $x\in X$ and the family of distributions is the admissibility density family
$\{p(\cdot\mid x): x\in X\}$.

The natural gradient, also due to Amari, is the gradient of a loss function $\ell(\theta)$
precondititioned by the inverse Fisher metric: $\tilde\nabla\ell(\theta) = g^{-1}(\theta)
\nabla\ell(\theta)$. It is the steepest ascent direction in the metric $g$, rather than
in the Euclidean metric on the parameter space. Learning algorithms that use the natural
gradient converge faster in information-geometric terms and are invariant under
reparameterization of the statistical manifold.

The $\alpha$-connections $\nabla^{(\alpha)}$ for $\alpha\in\mathbb{R}$ are a one-parameter
family of connections on the statistical manifold, of which the Levi-Civita connection
corresponds to $\alpha=0$ and the exponential and mixture connections correspond to
$\alpha=\pm 1$. The $\pm 1$-connections are flat (zero curvature) for exponential family
distributions, which is relevant for the analysis of admissibility manifolds arising from
constraint structures with exponential family structure.

\chapter{Differential Geometry: Curvature and Submersions}

A smooth submersion $\pi: X\to M$ between Riemannian manifolds is a smooth map whose
differential $d\pi_x: T_xX\to T_{\pi(x)}M$ is surjective at every $x\in X$. The
vertical bundle $V = \ker(d\pi)$ and its $g$-orthogonal complement, the horizontal bundle
$H = V^\perp$, provide the decomposition $TX = V\oplus H$ used throughout this monograph.

The sectional curvature $K_g(\sigma)$ of a Riemannian manifold $(X,g)$ at a point $x$ in
the direction of a two-plane $\sigma\subset T_xX$ is defined as:
\[
K_g(\sigma) = \frac{g(R(u,v)v,u)}{g(u,u)g(v,v) - g(u,v)^2}
\]
for any basis $\{u,v\}$ of $\sigma$, where $R$ is the Riemann curvature tensor
$R(u,v)w = \nabla_u\nabla_v w - \nabla_v\nabla_u w - \nabla_{[u,v]}w$. The projected
sectional curvature $\kpi(m)$ defined in Chapter~5 integrates $|K_g(v,h)|$ over mixed
planes $\sigma = \mathrm{span}\{v,h\}$ with $v\in V_x$ and $h\in H_x$.

\chapter{The O'Neill Formulas}

O'Neill's fundamental equations for a Riemannian submersion relate the curvatures of the
total space $(X,g)$, the base $(M, g_M)$, and the fibers to two tensor fields $A$ and $T$
called the O'Neill tensors~\cite{oneill1966}. Section~5.5 of the main text introduces
these tensors in connection with the collapse-driving curvature; this appendix provides
the complete formulas and their derivations.

\section*{The O'Neill Tensors: Full Definitions}

Let $\pi:(X,g)\to(M,g_M)$ be a Riemannian submersion with vertical bundle $V =
\ker(d\pi)$ and horizontal bundle $H = V^\perp$. For smooth vector fields $E, F$ on $X$,
write $E = E^H + E^V$ for the decomposition into horizontal and vertical parts. The
O'Neill $A$-tensor and $T$-tensor are defined as:
\begin{align*}
A_E F &= (\nabla_{E^H} F^V)^H + (\nabla_{E^H} F^H)^V, \\
T_E F &= (\nabla_{E^V} F^H)^V + (\nabla_{E^V} F^V)^H.
\end{align*}
The $A$-tensor is skew-symmetric in the sense that $g(A_E F, G) = -g(F, A_E G)$ for
horizontal $E$ and arbitrary $F, G$. The $T$-tensor at each fiber point is the second
fundamental form of the fiber in $X$.

\section*{The Complete Curvature Equations}

For orthonormal vectors $h_1, h_2 \in H_x$ (horizontal) and $v_1, v_2 \in V_x$
(vertical), the O'Neill curvature formulas are:
\begin{align}
K_X(h_1, h_2) &= K_M(d\pi(h_1), d\pi(h_2)) - 3\|A_{h_1}h_2\|^2, \label{eq:app-HH}\\
K_X(v, h) &= -g\!\left((\nabla_h A)_v h,\, v\right) - \|A_h v\|^2 + \|T_v h\|^2,
  \label{eq:app-VH}\\
K_X(v_1, v_2) &= K_F(v_1, v_2) - \|T_{v_1}v_2\|^2 + \|T_{v_1}v_1\|\cdot\|T_{v_2}v_2\| + \ldots,
  \label{eq:app-VV}
\end{align}
where $K_M$ is the sectional curvature of the base and $K_F$ is the sectional curvature
of the fiber. The horizontal-horizontal equation~\eqref{eq:app-HH} shows that the base
curvature always exceeds the total-space horizontal curvature by $3\|A_{h_1}h_2\|^2$:
the $A$-tensor reduces curvature as seen from the base.

The mixed equation~\eqref{eq:app-VH} is the key formula for the Projection-Collapse
Principle. Rearranging:
\[
K_X(v,h) = \|T_v h\|^2 - \|A_h v\|^2 - g\!\left((\nabla_h A)_v h,\, v\right).
\]
The collapse-driving term is $-\|A_h v\|^2$: the $A$-tensor contribution is negative,
meaning that non-integrability of the horizontal distribution \emph{reduces} the mixed
sectional curvature, coupling the vertical and horizontal directions. The projected
sectional curvature~\eqref{eq:kpi} is therefore controlled by $\|A\|^2$ as established
in equation~\eqref{eq:kpi-A}.

\section*{Integrability and Collapse}

The horizontal distribution $H$ is integrable (i.e., locally a product $V \times M$) if
and only if $A \equiv 0$. In this case all mixed sectional curvatures vanish, $\kpi = 0$,
and the Projection-Collapse Principle gives $\Lam = 0$: no collapse occurs. This is the
geometric condition for a representation to be perfectly faithful on the mixed directions.

When $A \not\equiv 0$, the horizontal distribution is non-integrable: horizontal curves
through adjacent fibers do not remain in the same fiber, and the fiber structure twists
relative to the base. This twisting is the geometric source of collapse, and $\|A\|^2$
is its quantitative measure.



\chapter{Proof Details for the Weak Form}

We provide a more detailed version of the proof of Theorem~\ref{thm:weak}. The key
measure-theoretic step is the following lemma.

\begin{lemma}\label{lem:divergence}
Let $\mu$ be a measure on a product space $\mathcal{Y}_1\times\mathcal{Y}_2$, and let
$\mu_1, \mu_2$ be its marginals. If there exists a set $E\subset\mathcal{Y}_1\times
\mathcal{Y}_2$ with $\mu(E)>0$ on which the density $d\mu/d(\mu_1\otimes\mu_2)$ differs
from 1, then $D(\mu, \mu_1\otimes\mu_2)>0$ for any information divergence $D$ that is
zero if and only if $\mu = \mu_1\otimes\mu_2$.
\end{lemma}

The proof of the weak form applies this lemma with $\mathcal{Y}_1\times\mathcal{Y}_2$
corresponding to the product of two admissibility classes in $\pi^{-1}(m)$, and with the
mixing in the pushforward distribution providing the set $E$ of positive measure where
the density differs from the product. The monotone function $f$ is constructed by
tracking the size of $E$ as a function of $\kpi(m)$ using the continuity of $\Lam$ and
the compactness of the fiber measure.

\chapter{Connection to RSVP Field Theory}

The Relativistic Scalar-Vector Plenum framework models the constraint structure of a
semantic domain as a coupled scalar-vector field over a spacetime manifold $(\mathcal{M},
\eta)$. The scalar field $\Phi:\mathcal{M}\to\mathbb{R}$ is the
admissibility density; the vector field $V:\mathcal{M}\to T\mathcal{M}$ is the
accessibility flow. The RSVP field equations are:
\begin{align}
\Box\Phi &= \lambda(\Phi^2 - v^2)\Phi - \nabla_\mu V^\mu, \\
\nabla_\nu F^{\mu\nu} &= J^\mu[\Phi],
\end{align}
where $\Box$ is the wave operator on $(\mathcal{M},\eta)$, $F^{\mu\nu}$ is the field
strength of $V$, and $J^\mu[\Phi]$ is the current sourced by the scalar field. These
equations enforce a Mexican-hat potential for $\Phi$, producing spontaneous symmetry
breaking at locations where the admissibility density falls below the threshold $v$: the
semantic bottlenecks discussed in Chapter~3.

The connection to the Fisher metric is through the identification $g_{ij}(x) = \partial_i
\partial_j \Phi(x)$ near local maxima of $\Phi$, where the scalar field is strictly
concave and the Hessian is negative definite. Near maxima, the RSVP scalar field
approximates the admissibility density, and its Hessian approximates the Fisher metric.
Away from maxima, the Fisher metric is the correct object. The RSVP framework provides
the dynamics; the Fisher metric provides the static geometry; they agree in the
high-admissibility regions where the system spends most of its time.

%% ── Bibliography ──────────────────────────────────────────────────────────
\backmatter

\begin{thebibliography}{99}

\bibitem{alemi2017}
A. Alemi, I. Fischer, J. Dillon, and K. Murphy,
``Deep variational information bottleneck,''
\textit{International Conference on Learning Representations}, 2017.

\bibitem{amari2016information}
S. Amari,
\textit{Information Geometry and Its Applications},
Springer, Tokyo, 2016.

\bibitem{amari2000methods}
S. Amari and H. Nagaoka,
\textit{Methods of Information Geometry},
American Mathematical Society, Providence, 2000.

\bibitem{bergson1911}
H. Bergson,
\textit{Creative Evolution},
translated by A. Mitchell, Henry Holt, New York, 1911.

\bibitem{bleecker1981}
D. Bleecker,
\textit{Gauge Theory and Variational Principles},
Addison-Wesley, Reading MA, 1981.

\bibitem{bohlen2026}
K. Bohlen,
``Reentrant value fields as delayed coupled reaction-diffusion systems on finite graphs,''
arXiv:2605.03940v2, 2026.

\bibitem{bredon1997}
G. E. Bredon,
\textit{Sheaf Theory}, second edition,
Springer, New York, 1997.

\bibitem{bronstein2021}
M. M. Bronstein, J. Bruna, T. Cohen, and P. Veli\v{c}kovi\'{c},
``Geometric deep learning: grids, groups, graphs, geodesics, and gauges,''
arXiv:2104.13478, 2021.

\bibitem{carlsson2009}
G. Carlsson,
``Topology and data,''
\textit{Bulletin of the American Mathematical Society} 46 (2009), 255--308.

\bibitem{cencov1982}
N. N. Chentsov,
\textit{Statistical Decision Rules and Optimal Inference},
American Mathematical Society, Providence, 1982.

\bibitem{cheeger1975}
J. Cheeger and D. Ebin,
\textit{Comparison Theorems in Riemannian Geometry},
North-Holland, Amsterdam, 1975.

\bibitem{clark2016}
A. Clark,
\textit{Surfing Uncertainty: Prediction, Action, and the Embodied Mind},
Oxford University Press, New York, 2016.

\bibitem{denham2025}
S. Denham,
``Semantic collapse in embedding space: neighborhood semantics and entropic drift,''
2026.

\bibitem{edelsbrunner2010}
H. Edelsbrunner and J. Harer,
\textit{Computational Topology: An Introduction},
American Mathematical Society, Providence, 2010.

\bibitem{friston2010}
K. Friston,
``The free-energy principle: a unified brain theory?''
\textit{Nature Reviews Neuroscience} 11 (2010), 127--138.

\bibitem{kaplan1989}
D. Kaplan,
``Demonstratives,''
in J. Almog, J. Perry, and H. Wettstein (eds.),
\textit{Themes from Kaplan},
Oxford University Press, Oxford, 1989, pp.~481--563.

\bibitem{kripke1963}
S. A. Kripke,
``Semantical analysis of modal logic I: normal modal propositional calculi,''
\textit{Zeitschrift f\"ur Mathematische Logik und Grundlagen der Mathematik} 9 (1963),
67--96.

\bibitem{lee2018}
J. M. Lee,
\textit{Introduction to Riemannian Manifolds}, second edition,
Springer, Cham, 2018.

\bibitem{marsden1999}
J. E. Marsden and T. S. Ratiu,
\textit{Introduction to Mechanics and Symmetry}, second edition,
Springer, New York, 1999.

\bibitem{ollivier2009}
Y. Ollivier,
``Ricci curvature of Markov chains on metric spaces,''
\textit{Journal of Functional Analysis} 256 (2009), 810--864.

\bibitem{oneill1966}
B. O'Neill,
``The fundamental equations of a submersion,''
\textit{Michigan Mathematical Journal} 13 (1966), 459--469.

\bibitem{oneill1983}
B. O'Neill,
\textit{Semi-Riemannian Geometry with Applications to Relativity},
Academic Press, New York, 1983.

\bibitem{perry1993}
J. Perry,
\textit{The Problem of the Essential Indexical and Other Essays},
Oxford University Press, New York, 1993.

\bibitem{petersen2016}
P. Petersen,
\textit{Riemannian Geometry}, third edition,
Springer, Cham, 2016.

\bibitem{robinson2014}
M. Robinson,
\textit{Topological Signal Processing},
Springer, Heidelberg, 2014.

\bibitem{searle1980}
J. R. Searle,
``Minds, brains, and programs,''
\textit{Behavioral and Brain Sciences} 3 (1980), 417--457.

\bibitem{shannon1948}
C. E. Shannon,
``A mathematical theory of communication,''
\textit{Bell System Technical Journal} 27 (1948), 379--423.

\bibitem{tikhonov1977}
A. N. Tikhonov and V. Y. Arsenin,
\textit{Solutions of Ill-Posed Problems},
Winston, Washington DC, 1977.

\bibitem{tishby2000}
N. Tishby, F. Pereira, and W. Bialek,
``The information bottleneck method,''
arXiv:physics/0004057, 2000.

\bibitem{villani2003}
C. Villani,
\textit{Topics in Optimal Transportation},
American Mathematical Society, Providence, 2003.

\bibitem{villani2009}
C. Villani,
\textit{Optimal Transport: Old and New},
Springer, Berlin, 2009.

\bibitem{whitehead1929}
A. N. Whitehead,
\textit{Process and Reality},
Macmillan, New York, 1929.

\end{thebibliography}

\end{document}
